% Generated by the blueprint tool (blueprint-pdf/1) from a blueprint snapshot; edit the
% Markdown sources, not this file.  Compile with `latexmk -pdf` or by running
% `pdflatex` twice (the second run settles the contents and the page references).
% Images, if any, are read from assets/ next to this file.
\documentclass[11pt]{article}
\usepackage[T1]{fontenc}
\usepackage[utf8]{inputenc}
\usepackage{lmodern}
\usepackage{textcomp}
\usepackage[a4paper,margin=2.5cm,headheight=14pt]{geometry}
\usepackage{amsmath,amssymb,amsthm}
\usepackage{amscd}  % KaTeX's CD diagrams; part of amsmath, so always there
\usepackage{mathtools}
\usepackage{bm}
\usepackage{graphicx}
\usepackage{xcolor}
\usepackage{array,booktabs,tabularx}
\usepackage{fancyvrb}
\usepackage[normalem]{ulem}
\usepackage{microtype}
% Optional, and only with their fonts: a small TeX installation has the
% packages without the fonts, which fails at the first formula.
\IfFileExists{mathrsfs.sty}{\IfFileExists{rsfs10.tfm}{\usepackage{mathrsfs}}{}}{}
\IfFileExists{stmaryrd.sty}{\IfFileExists{stmary10.tfm}{\usepackage{stmaryrd}}{}}{}
\IfFileExists{cancel.sty}{\usepackage{cancel}}{}
\usepackage[hyphens]{url}
\usepackage[pdfusetitle,bookmarksnumbered=false]{hyperref}
\usepackage{bookmark}
\definecolor{bplink}{HTML}{1F4E9C}
\definecolor{bpmuted}{HTML}{5F676E}
\hypersetup{colorlinks=true,linkcolor=bplink,urlcolor=bplink,citecolor=bplink}
\setlength{\parindent}{0pt}
\setlength{\parskip}{0.55em plus 0.2em minus 0.1em}
% Lists as tight as the page's.
\makeatletter
\def\@listi{\leftmargin\leftmargini\parsep 0pt\topsep 0.3em plus 0.1em\itemsep 0.1em plus 0.1em}
\let\@listI\@listi
\def\@listii{\leftmargin\leftmarginii\labelwidth\leftmarginii\advance\labelwidth-\labelsep\parsep 0pt\topsep 0.1em\itemsep 0.1em}
% An ordered list that does not start at 1.
\newcommand{\bpstartat}[1]{\setcounter{\@enumctr}{\numexpr#1-1\relax}}
\makeatother
\setcounter{tocdepth}{3}
\emergencystretch=3em
\allowdisplaybreaks

% A numbered statement, set as a paper sets one: "Theorem 2.1 (Title)."
% #1 is the lead ("Theorem~2.1"), #2 what follows it in roman.
\newenvironment{bpstatement}[2]{\par\addvspace{0.6\baselineskip}\noindent\textbf{#1}#2\textbf{.}\enspace\ignorespaces}
  {\par\addvspace{0.4\baselineskip}}
\newenvironment{bpstatementit}[2]{\begin{bpstatement}{#1}{#2}\itshape\ignorespaces}{\end{bpstatement}}
% A step: a relation with prose of its own, read under the object it starts from.
\newenvironment{bpstep}[1]{\par\addvspace{0.3\baselineskip}\begin{list}{}{\leftmargin=1.6em\rightmargin=0pt\topsep=0pt}\item\relax
  {\small\color{bpmuted}#1}\par\nopagebreak}{\end{list}\par}
% Notes about where an object is written out.
\newenvironment{bpnote}{\par\small\color{bpmuted}\itshape}{\par}
\newcommand{\bpomitted}[1]{\par{\small\color{bpmuted}\itshape #1\par}}
\newcommand{\bpbodyheading}[1]{\par\addvspace{0.4\baselineskip}\noindent\textbf{#1}\par\nopagebreak}
\newcommand{\bprule}{\par\medskip\noindent\hfil\rule{0.3\linewidth}{0.4pt}\hfil\par\medskip}
\newcommand{\bpboundary}[1]{\par\nopagebreak{\small\color{bpmuted}#1\par}\nopagebreak}
% A reference to something this export leaves out, and one that resolves nowhere.
\newcommand{\bpout}[1]{#1\textsuperscript{\hyperlink{bp-outside}{\ensuremath{\ast}}}}
\newcommand{\bpbroken}[1]{\textcolor{red!70!black}{[#1?]}}
\newcommand{\bpslug}[1]{\texttt{#1}}
% A bibliography entry: label, anchor, entry; the label hangs in the margin.
\newcommand{\bpbibitem}[3]{\par\setbox0\hbox{[#1]\enspace}\ifdim\wd0<5em\setbox0\hbox to 5em{[#1]\hfil}\fi
  \hangindent=\wd0 \hangafter=1 \noindent\box0\phantomsection\label{#2}#3\par\smallskip}
% A character the export knows no way to set.
\newcommand{\bpglyph}[1]{\fbox{\scriptsize #1}}
% Images, at their own size unless they are wider than the line.
\newsavebox{\bpimagebox}
\newcommand{\bpimage}[1]{\sbox{\bpimagebox}{\includegraphics{#1}}%
  \ifdim\wd\bpimagebox>\linewidth\resizebox{\linewidth}{!}{\usebox{\bpimagebox}}\else\usebox{\bpimagebox}\fi}
\newcommand{\bpmissingimage}[1]{\fbox{\parbox{0.8\linewidth}{\centering\small\color{bpmuted}Image not included: #1}}}
\DefineVerbatimEnvironment{bpverbatim}{Verbatim}{fontsize=\small,xleftmargin=1em}
\renewcommand{\qedsymbol}{\ensuremath{\blacksquare}}
\hypersetup{pdftitle={Weil Conjectures — E Specialization and contracting correspondences},pdfsubject={Appendix E Specialization and contracting correspondences; all levels},pdfcreator={blueprint}}
% KaTeX commands the website knows
\providecommand{\R}{\mathbb{R}}\providecommand{\N}{\mathbb{N}}\providecommand{\Z}{\mathbb{Z}}
\providecommand{\Q}{\mathbb{Q}}\providecommand{\C}{\mathbb{C}}
\providecommand{\Reals}{\mathbb{R}}\providecommand{\reals}{\mathbb{R}}\providecommand{\Complex}{\mathbb{C}}
\providecommand{\cnums}{\mathbb{C}}\providecommand{\natnums}{\mathbb{N}}\providecommand{\Bbb}[1]{\mathbb{#1}}
\providecommand{\bold}[1]{\mathbf{#1}}\providecommand{\frak}[1]{\mathfrak{#1}}\providecommand{\mathbbm}[1]{\mathbb{#1}}
\providecommand{\lang}{\langle}\providecommand{\rang}{\rangle}\providecommand{\bra}[1]{\langle #1\rvert}
\providecommand{\ket}[1]{\lvert #1\rangle}\providecommand{\braket}[1]{\langle #1\rangle}
\providecommand{\set}[1]{\{#1\}}\providecommand{\Set}[1]{\left\{#1\right\}}
\providecommand{\infin}{\infty}\providecommand{\isin}{\in}\providecommand{\plusmn}{\pm}\providecommand{\sub}{\subset}
\providecommand{\sube}{\subseteq}\providecommand{\supe}{\supseteq}\providecommand{\exist}{\exists}
\providecommand{\bull}{\bullet}\providecommand{\larr}{\leftarrow}\providecommand{\rarr}{\rightarrow}
\providecommand{\lrarr}{\leftrightarrow}\providecommand{\Larr}{\Leftarrow}\providecommand{\Rarr}{\Rightarrow}
\providecommand{\Lrarr}{\Leftrightarrow}\providecommand{\uarr}{\uparrow}\providecommand{\darr}{\downarrow}
\providecommand{\harr}{\leftrightarrow}\providecommand{\lArr}{\Leftarrow}\providecommand{\rArr}{\Rightarrow}
\providecommand{\hArr}{\Leftrightarrow}\providecommand{\alef}{\aleph}\providecommand{\alefsym}{\aleph}
\providecommand{\real}{\Re}\providecommand{\image}{\Im}\providecommand{\weierp}{\wp}
\providecommand{\Alpha}{\mathrm{A}}\providecommand{\Beta}{\mathrm{B}}\providecommand{\Epsilon}{\mathrm{E}}
\providecommand{\Zeta}{\mathrm{Z}}\providecommand{\Eta}{\mathrm{H}}\providecommand{\Iota}{\mathrm{I}}
\providecommand{\Kappa}{\mathrm{K}}\providecommand{\Mu}{\mathrm{M}}\providecommand{\Nu}{\mathrm{N}}
\providecommand{\Omicron}{\mathrm{O}}\providecommand{\Rho}{\mathrm{P}}\providecommand{\Tau}{\mathrm{T}}
\providecommand{\Chi}{\mathrm{X}}\providecommand{\omicron}{o}\providecommand{\gt}{>}\providecommand{\lt}{<}
\providecommand{\And}{\&}\providecommand{\hdashline}{\hline}\providecommand{\operatornamewithlimits}[1]{\operatorname*{#1}}
\providecommand{\llbracket}{[\![}\providecommand{\rrbracket}{]\!]}\providecommand{\cancel}[1]{#1}
\providecommand{\mathscr}{\mathcal}
% the characters beyond ASCII this document uses
\DeclareUnicodeCharacter{2160}{\textup{I}}
\DeclareUnicodeCharacter{2161}{\textup{II}}
% the project's macros ([katex.macros] in blueprint.toml)
\providecommand{\A}{}\renewcommand{\A}{\mathbf A}
\providecommand{\AEn}{}\renewcommand{\AEn}{A_{E,n}}
\providecommand{\Aaf}{}\renewcommand{\Aaf}{\mathbf A}
\providecommand{\Ab}{}\renewcommand{\Ab}{\mathsf{Ab}}
\providecommand{\AffSch}{}\renewcommand{\AffSch}{\operatorname{AffSch}}
\providecommand{\Aone}{}\renewcommand{\Aone}{\mathbf A^1}
\providecommand{\Bl}{}\renewcommand{\Bl}{\operatorname{Bl}}
\providecommand{\Br}{}\renewcommand{\Br}{\mathrm{Br}}
\providecommand{\C}{}\renewcommand{\C}{\mathbb C}
\providecommand{\Cech}{}\renewcommand{\Cech}{\check{H}}
\providecommand{\Coext}{}\renewcommand{\Coext}{\operatorname{Coext}}
\providecommand{\Cone}{}\renewcommand{\Cone}{\operatorname{Cone}}
\providecommand{\Cons}{}\renewcommand{\Cons}{\operatorname{Cons}}
\providecommand{\Cref}{}\renewcommand{\Cref}[1]{\text{#1}}
\providecommand{\D}{}\renewcommand{\D}{D}
\providecommand{\DLAFT}{}\renewcommand{\DLAFT}{\operatorname{FT}}
\providecommand{\DLAFr}{}\renewcommand{\DLAFr}{\operatorname{Fr}}
\providecommand{\DLAQ}{}\renewcommand{\DLAQ}{{\overline{\mathbf Q}_\ell}}
\providecommand{\DLRQ}{}\renewcommand{\DLRQ}{{\overline{\mathbf Q}_\ell}}
\providecommand{\Dbc}{}\renewcommand{\Dbc}{D^b_c}
\providecommand{\Dc}{}\renewcommand{\Dc}{D_{\mathrm{cons}}}
\providecommand{\Dctf}{}\renewcommand{\Dctf}{D^b_{ctf}}
\providecommand{\Dic}{}\renewcommand{\Dic}{D_{\IndCons}}
\providecommand{\Dplus}{}\renewcommand{\Dplus}{D^{+}}
\providecommand{\End}{}\renewcommand{\End}{\operatorname{End}}
\providecommand{\Et}{}\renewcommand{\Et}{\mathrm{\acute{e}t}}
\providecommand{\Ex}{}\renewcommand{\Ex}{\mathrm{Ex}}
\providecommand{\Ext}{}\renewcommand{\Ext}{\operatorname{Ext}}
\providecommand{\F}{}\renewcommand{\F}{\mathbb F}
\providecommand{\FT}{}\renewcommand{\FT}{\operatorname{FT}}
\providecommand{\Fbar}{}\renewcommand{\Fbar}{{\overline{\mathbf F}_q}}
\providecommand{\Fell}{}\renewcommand{\Fell}{\mathbf F_\ell}
\providecommand{\Fib}{}\renewcommand{\Fib}{\operatorname{Fib}}
\providecommand{\Fix}{}\renewcommand{\Fix}{\operatorname{Fix}}
\providecommand{\Fp}{}\renewcommand{\Fp}{{\mathbf F_{p}}}
\providecommand{\Fq}{}\renewcommand{\Fq}{{\mathbf F_q}}
\providecommand{\Fqn}{}\renewcommand{\Fqn}{{\mathbf F_{q^{n}}}}
\providecommand{\Frob}{}\renewcommand{\Frob}{\operatorname{Frob}}
\providecommand{\GL}{}\renewcommand{\GL}{\operatorname{GL}}
\providecommand{\Gal}{}\renewcommand{\Gal}{\operatorname{Gal}}
\providecommand{\Hens}{}\renewcommand{\Hens}{\operatorname{Hens}}
\providecommand{\Hilb}{}\renewcommand{\Hilb}{\operatorname{Hilb}}
\providecommand{\Hom}{}\renewcommand{\Hom}{\operatorname{Hom}}
\providecommand{\IndCons}{}\renewcommand{\IndCons}{\operatorname{IndCons}}
\providecommand{\Kar}{}\renewcommand{\Kar}{\operatorname{Kar}}
\providecommand{\LT}{}\renewcommand{\LT}{\operatorname{LT}}
\providecommand{\Lie}{}\renewcommand{\Lie}{\operatorname{Lie}}
\providecommand{\Loc}{}\renewcommand{\Loc}{\operatorname{Loc}}
\providecommand{\Mod}{}\renewcommand{\Mod}{\mathsf{Mod}}
\providecommand{\N}{}\renewcommand{\N}{\mathbb{N}}
\providecommand{\Needspace}{}\renewcommand{\Needspace}{\par}
\providecommand{\O}{}\renewcommand{\O}{{\mathcal{O}}}
\providecommand{\Ob}{}\renewcommand{\Ob}{\mathrm{Ob}}
\providecommand{\Oe}{}\renewcommand{\Oe}{\mathcal O_E}
\providecommand{\PP}{}\renewcommand{\PP}{\mathbf{P}}
\providecommand{\Perf}{}\renewcommand{\Perf}{\operatorname{Perf}}
\providecommand{\Pic}{}\renewcommand{\Pic}{\mathrm{Pic}}
\providecommand{\Pm}{}\renewcommand{\Pm}{\mathbf P^{m}}
\providecommand{\Pone}{}\renewcommand{\Pone}{\mathbf P^{1}}
\providecommand{\Pqn}{}\renewcommand{\Pqn}{\mathbf P^{n}}
\providecommand{\ProFin}{}\renewcommand{\ProFin}{\operatorname{ProFin}}
\providecommand{\Proet}{}\renewcommand{\Proet}{\mathrm{pro\acute{e}t}}
\providecommand{\Q}{}\renewcommand{\Q}{\mathbb Q}
\providecommand{\Qb}{}\renewcommand{\Qb}{\overline{\mathbf Q}_{\ell}}
\providecommand{\Qbar}{}\renewcommand{\Qbar}{\overline{\mathbf Q}}
\providecommand{\QbarX}{}\renewcommand{\QbarX}{\overline{\mathbf Q}_{\ell,X}}
\providecommand{\Qbarl}{}\renewcommand{\Qbarl}{{\overline{\mathbf Q}_{\ell}}}
\providecommand{\Qell}{}\renewcommand{\Qell}{{\mathbf{Q}_{\ell}}}
\providecommand{\Ql}{}\renewcommand{\Ql}{\mathbf Q_{\ell}}
\providecommand{\Qlb}{}\renewcommand{\Qlb}{{\overline{\mathbf Q}_\ell}}
\providecommand{\Qlbar}{}\renewcommand{\Qlbar}{\overline{\mathbf Q}_\ell}
\providecommand{\R}{}\renewcommand{\R}{\mathbb{R}}
\providecommand{\RG}{}\renewcommand{\RG}{\operatorname{\mathrm{R}\Gamma}}
\providecommand{\RGamma}{}\renewcommand{\RGamma}{R\Gamma}
\providecommand{\RGammac}{}\renewcommand{\RGammac}{R\Gamma_c}
\providecommand{\RHom}{}\renewcommand{\RHom}{R\mathcal{H}om}
\providecommand{\RHomX}{}\renewcommand{\RHomX}{R\mathcal H om_X}
\providecommand{\Rat}{}\renewcommand{\Rat}{\operatorname{IsRat}}
\providecommand{\Res}{}\renewcommand{\Res}{\operatorname{Res}}
\providecommand{\Rlim}{}\renewcommand{\Rlim}{R\!\varprojlim}
\providecommand{\SIm}{}\renewcommand{\SIm}{\operatorname{SIm}}
\providecommand{\Sch}{}\renewcommand{\Sch}{\mathsf{Sch}}
\providecommand{\Sh}{}\renewcommand{\Sh}{\operatorname{Sh}}
\providecommand{\Span}{}\renewcommand{\Span}{\operatorname{span}}
\providecommand{\Spec}{}\renewcommand{\Spec}{\operatorname{Spec}}
\providecommand{\Tr}{}\renewcommand{\Tr}{\operatorname{Tr}}
\providecommand{\Trsm}{}\renewcommand{\Trsm}{\operatorname{Tr}^{\mathrm{sm}}}
\providecommand{\Wt}{}\renewcommand{\Wt}{\operatorname{Wt}}
\providecommand{\Z}{}\renewcommand{\Z}{\mathbb{Z}}
\providecommand{\Zar}{}\renewcommand{\Zar}{\mathrm{zar}}
\providecommand{\Zell}{}\renewcommand{\Zell}{{\mathbf{Z}_{\ell}}}
\providecommand{\Zl}{}\renewcommand{\Zl}{{\mathbf Z_\ell}}
\providecommand{\affet}{}\renewcommand{\affet}[1]{{#1}^{\mathrm{aff}}_{\mathrm{ét}}}
\providecommand{\affproet}{}\renewcommand{\affproet}[1]{{#1}^{\mathrm{aff}}_{\mathrm{proét}}}
\providecommand{\boxt}{}\renewcommand{\boxt}{\boxtimes_A^{\mathrm L}}
\providecommand{\cF}{}\renewcommand{\cF}{\mathcal{F}}
\providecommand{\cHom}{}\renewcommand{\cHom}{R\mathcal{H}\!om}
\providecommand{\cI}{}\renewcommand{\cI}{\mathcal{I}}
\providecommand{\cL}{}\renewcommand{\cL}{\mathcal L}
\providecommand{\cM}{}\renewcommand{\cM}{\mathcal M}
\providecommand{\cN}{}\renewcommand{\cN}{\mathcal N}
\providecommand{\cO}{}\renewcommand{\cO}{\mathcal O}
\providecommand{\cT}{}\renewcommand{\cT}{\mathcal T}
\providecommand{\checked}{}\renewcommand{\checked}{\color{red}\aftergroup\restorecolor}
\providecommand{\coeff}{}\renewcommand{\coeff}{\operatorname{coeff}}
\providecommand{\coker}{}\renewcommand{\coker}{\operatorname{coker}}
\providecommand{\colim}{}\renewcommand{\colim}{\operatorname{colim}}
\providecommand{\constructedby}{}\renewcommand{\constructedby}[1]{\par\smallskip\noindent\textcolor{purple!60!black}{\textsf{Constructed by:} }
  \footnotesize\sloppy\nolinkurl{#1}\par}
\providecommand{\deepresult}{}\renewcommand{\deepresult}[1]{\par\smallskip\noindent\textcolor{orange!45!black}{
  \textsf{Blueprint role: named deep theorem node; see the cluster proof architecture.}}
  \par\noindent\textcolor{orange!45!black}{\textsf{Mathematical source:} }#1\par}
\providecommand{\dt}{}\renewcommand{\dt}{\mathbin{\otimes^{\mathbf L}}}
\providecommand{\dw}{}\renewcommand{\dw}{\operatorname{dw}}
\providecommand{\et}{}\renewcommand{\et}[1]{{#1}_{\mathrm{ét}}}
\providecommand{\ev}{}\renewcommand{\ev}{\operatorname{ev}}
\providecommand{\fib}{}\renewcommand{\fib}{\operatorname{fib}}
\providecommand{\formalizationnote}{}\renewcommand{\formalizationnote}[1]{\begin{quote}\small\raggedright\textbf{Formalization note.} #1\end{quote}}
\providecommand{\fpqc}{}\renewcommand{\fpqc}{X_{\mathrm{fpqc}}}
\providecommand{\gammacont}{}\renewcommand{\gammacont}{\Gamma_{\mathrm{cont}}}
\providecommand{\gm}{}\renewcommand{\gm}{\mathbf G_m}
\providecommand{\hocolim}{}\renewcommand{\hocolim}{\operatorname{hocolim}}
\providecommand{\holim}{}\renewcommand{\holim}{\operatorname{holim}}
\providecommand{\hrB}{}\renewcommand{\hrB}{\@ifnextchar]{\hrEnd}{\hrStep}}
\providecommand{\hrEnd}{}\renewcommand{\hrEnd}{\@firstoftwo{}}
\providecommand{\hrStep}{}\renewcommand{\hrStep}{\@firstoftwo{\hrB}}
\providecommand{\id}{}\renewcommand{\id}{\operatorname{id}}
\providecommand{\kunFl}{}\renewcommand{\kunFl}{\mathbf F_\ell}
\providecommand{\kunZn}{}\renewcommand{\kunZn}{\mathbf Z/n\mathbf Z}
\providecommand{\kuncd}{}\renewcommand{\kuncd}{\operatorname{cd}}
\providecommand{\kundim}{}\renewcommand{\kundim}{\dim}
\providecommand{\kunotimes}{}\renewcommand{\kunotimes}[1]{\otimes^{\mathbf L}_{#1}}
\providecommand{\kunsupp}{}\renewcommand{\kunsupp}{\operatorname{Supp}}
\providecommand{\lotimes}{}\renewcommand{\lotimes}{\otimes^{\mathrm L}}
\providecommand{\m}{}\renewcommand{\m}{{\mathfrak{m}}}
\providecommand{\n}{}\renewcommand{\n}{{\mathfrak{n}}}
\providecommand{\nagataGap}{}\renewcommand{\nagataGap}[1]{\par\noindent\textbf{Formalization gap.} #1\par}
\providecommand{\openobligation}{}\renewcommand{\openobligation}[1]{\par\smallskip\noindent\textcolor{red!65!black}{\textsf{Open obligation:} }
  \footnotesize\sloppy\nolinkurl{#1}\par}
\providecommand{\p}{}\renewcommand{\p}{{\mathfrak{p}}}
\providecommand{\pr}{}\renewcommand{\pr}{\operatorname{pr}}
\providecommand{\preshv}{}\renewcommand{\preshv}[1]{\mathrm{PreShv}(#1)}
\providecommand{\pro}{}\renewcommand{\pro}[1]{\mathrm{Pro}(#1)}
\providecommand{\proet}{}\renewcommand{\proet}{X_{\mathrm{proét}}}
\providecommand{\q}{}\renewcommand{\q}{{\mathfrak{q}}}
\providecommand{\rad}{}\renewcommand{\rad}{\operatorname{rad}}
\providecommand{\rank}{}\renewcommand{\rank}{\operatorname{rank}}
\providecommand{\red}{}\renewcommand{\red}{\color{red}}
\providecommand{\restorecolor}{}\renewcommand{\restorecolor}{\color{black}}
\providecommand{\retiredobligation}{}\renewcommand{\retiredobligation}[1]{\par\smallskip\noindent\textcolor{red!75!black}{\textsf{Retired invalid route:} }
  \footnotesize\sloppy\nolinkurl{#1}\par}
\providecommand{\rk}{}\renewcommand{\rk}{\operatorname{rk}}
\providecommand{\sep}{}\renewcommand{\sep}{\textnormal{sep}}
\providecommand{\shv}{}\renewcommand{\shv}[1]{\mathrm{Shv}(#1)}
\providecommand{\spec}{}\renewcommand{\spec}[1]{\mathrm{Spec}(#1)}
\providecommand{\stacks}{}\renewcommand{\stacks}[1]{\href{https://stacks.math.columbia.edu/tag/#1}{\textnormal{Stacks, Tag~#1}}}
\providecommand{\stacksproject}{}\renewcommand{\stacksproject}[1]{{\cite[\href{https://stacks.math.columbia.edu/tag/#1}{Tag #1}]{stacks-project}}}
\providecommand{\statusnote}{}\renewcommand{\statusnote}[1]{\begin{quote}\small\raggedright\textbf{Status.} #1\end{quote}}
\providecommand{\tildeideal}{}\renewcommand{\tildeideal}[2]{\widetilde{(#1, #2)}}
\providecommand{\tildering}{}\renewcommand{\tildering}[3]{{#1}_{\widetilde{(#2, #3)}}}
\providecommand{\tr}{}\renewcommand{\tr}{\operatorname{tr}}
\providecommand{\wt}{}\renewcommand{\wt}{\operatorname{wt}}

\begin{document}
\pagestyle{plain}
\begin{center}
{\LARGE\bfseries Weil Conjectures\par}
\vspace{0.7em}
{\large Appendix E: Specialization and contracting correspondences\par}
\vspace{0.5em}
{\small\color{bpmuted}Snapshot 2026-09-26, revision 8e4cb88e2a\par}
\end{center}
\vspace{0.5em}
{\small\tableofcontents}
\bigskip\hrule\bigskip

% object: chap-Specialization-and-contracting-correspondences (blueprint/chap-Specialization-and-contracting-correspondences/_section.md)
\section*{\texorpdfstring{E\quad Specialization and contracting correspondences}{E Specialization and contracting correspondences}}
\phantomsection\label{bp1}%
\addcontentsline{toc}{section}{\texorpdfstring{E\quad Specialization and contracting correspondences}{E Specialization and contracting correspondences}}
% text: prose-E-01
\phantomsection\label{bp2}%
The trace formula needs the local term at an isolated fixed point to be the trace on the stalk. This is not a formal fact. A local term is computed from \(\cHom(c_2^*K,c_1^!K)\) near the fixed point, and the stalk does not determine this complex (Remark \bpout{12.4}). Instead, Chapter \bpout{12} proves the equality by a computation for Frobenius on a smooth curve (Proposition \bpout{12.16}) and reduces everything else to that case. In particular, its trace formula (Theorem \bpout{12.24}) does not depend on this appendix.

The computation leaves open why Frobenius behaves so well. Topology suggests that the reason is contraction. If a map \(\varphi\) of a manifold contracts a neighbourhood of a fixed point \(y\) into itself, then the local index at \(y\) is \(1\). With coefficients in a sheaf, the local index is the trace of \(\varphi\) on the stalk at \(y\). Frobenius contracts in the strongest possible sense: its differential is zero. In this appendix we make this explanation into a theorem, in any dimension and for a whole class of correspondences at once.

In algebraic geometry we measure contraction by powers of ideals. Let \(c=(c_1,c_2):C\to X\times X\) be a correspondence, \(Z\subset X\) a closed subscheme with ideal \(I\), and put \(I_i=c_i^{-1}(I)\mathcal O_C\). We say, as in [\hyperref[bp98]{Var05}], that \(c\) contracts towards \(Z\) if \(I_2\subset I_1\) and \(I_2^n\subset I_1^{n+1}\) for some \(n\geq1\). Frobenius is the model case. Indeed, take \(c=(\mathrm{id},\Phi_X^m)\) and \(Z\) a point \(y\in X_0(\mathbf F_{q^m})\), and choose coordinates \(x_i\) of an affine neighbourhood of \(y\) defined over \(\mathbf F_q\). Then \(y\) has coordinates \(a_i\in\mathbf F_{q^m}\), and \((\Phi_X^m)^*(x_i-a_i)=x_i^{q^m}-a_i=(x_i-a_i)^{q^m}\). Hence \(I_2\subset I_1^{q^m}\), and \(n=1\) works.

The main theorem, Theorem \hyperref[bp91]{E.30}, says that near a component of its fixed locus that meets \(Z\), the trace class of a contracting correspondence can be computed on \(Z\) alone, from the correspondence restricted to \(Z\). For \(Z\) a point, Corollary \hyperref[bp95]{E.31} says that the local term at an isolated fixed point is the trace on the stalk, even when \(X\) is singular. The theorem and its proof are due to Varshavsky [\hyperref[bp98]{Var05}]. They generalize Fujiwara's theorem, conjectured by Deligne, that the local terms of a correspondence composed with a sufficiently high power of Frobenius are the naive ones.

The proof makes the contraction literal by deforming \(X\) to the normal cone of \(Z\). Over the trait \(\operatorname{Spec}k[t]_{(t)}\), the deformation to the normal cone \(\mathcal D_ZX\) has generic fibre \(X\) and special fibre the normal cone \(N_ZX\). On the special fibre, the contraction hypothesis says that the second leg maps everything into the zero section (Lemma \hyperref[bp82]{E.26}). There the part of \(K\) away from \(Z\) contributes nothing. What remains is the specialization of \(K\) restricted to the zero section, and by Verdier's theorem (Theorem \hyperref[bp24]{E.8}) this is \(K|_Z\).

It remains to transport the answer back to \(X\). We identify the trace class on the special fibre with the specialization of the trace class of \(c\) (Proposition \hyperref[bp69]{E.22}). The specialization map is an isomorphism, because the fixed locus does not move in the family up to nilpotents (Lemma \hyperref[bp86]{E.28}). Before any of this, we must be able to specialize a trace class at all. A trace class is built from \(K\boxtimes\mathbb DK\), so we need nearby cycles to commute with external products and with Verdier duality. Propositions \hyperref[bp28]{E.9} and \hyperref[bp56]{Proposition E.18} prove both, by a dimension argument of Illusie [\hyperref[bp99]{Ill94}].

We use the six operations and the notation for correspondences and trace classes of Chapter \bpout{11}, and classical results on nearby cycles over a trait, which we recall where they are first needed. Throughout this appendix, \(k\) is algebraically closed, \(\ell\) is invertible in \(k\), and \(A\) is a finite commutative ring annihilated by a power of \(\ell\). The only exception to these conventions is the base of nearby cycles in Sections \hyperref[bp26]{E.2}–\hyperref[bp67]{Section E.5}, which may be any excellent strictly henselian trait, with residue field only separably closed. As in the rest of the notes, statements about constructible complexes on \(k\)-schemes are made on the pro-étale site, but we construct nearby cycles on the classical étale site and transport them by the comparison for constructible complexes (Theorem \bpout{2.43} and Corollary \bpout{2.44}).

% object: sec-Specialization-to-the-normal-cone (blueprint/chap-Specialization-and-contracting-correspondences/sec-Specialization-to-the-normal-cone/_section.md), content hidden
\subsection*{\texorpdfstring{E.1\quad Specialization to the normal cone}{E.1 Specialization to the normal cone}}
\phantomsection\label{bp3}%
\addcontentsline{toc}{subsection}{\texorpdfstring{E.1\quad Specialization to the normal cone}{E.1 Specialization to the normal cone}}
% text: prose-E-02
\phantomsection\label{bp4}%
Let \(S\) be the strict henselization of \(\operatorname{Spec}k[t]_{(t)}\), with generic point \(\eta\). Nearby cycles move a complex from the generic fibre of a family over \(S\) to its special fibre. Since the special fibre does not lie in the generic fibre, we cannot simply restrict. The obvious attempt is to extend by \(Rj_*\) along the open generic fibre and restrict by \(i^*\) to the special fibre. This fails already for \(S\) itself, because the fraction field of \(S\) is not separably closed. Its absolute Galois group is the inertia group, whose tame quotient is \(\prod_{\ell'\neq p}\mathbf Z_{\ell'}(1)\). For \(n\) invertible in \(k\), the stalk of \(R^1j_*(\mathbf Z/n)\) at the closed point is therefore \(H^1(\eta,\mathbf Z/n)=\mathbf Z/n(-1)\neq0\). This class comes from the arithmetic of \(\eta\), not from the geometry of the family. Passing first to the geometric generic fibre removes such Galois classes. After this change, the constant sheaf specializes to itself, in degree zero, by the first calculation below.

Both specialization to the normal cone and the comparison of trace classes in Proposition \hyperref[bp69]{E.22} are built from nearby cycles on the deformation to the normal cone. The key input is Verdier's theorem, Theorem \hyperref[bp24]{E.8} (see [\hyperref[bp98]{Var05}, §§1.3–1.4 and 3]).

The basic computation is that nearby cycles of a constant family \(Y\times S\) give back the complex we started with (Lemma \hyperref[bp7]{E.2}). We use it twice. It makes the map \(\tau_{X,Z}\) below land in \(i^*L\), and it makes the specialization of trace classes the identity when the fixed scheme does not move (Corollary \hyperref[bp78]{E.25}). Its proof uses only the definition of nearby cycles, their base change maps and two calculations, which we now recall.

Choose a geometric generic point \(\bar\eta\) of \(S\), and let \(\mathcal Y\to S\) be separated of finite type. To define nearby cycles, we pull \(\mathcal Y\) back to the integral closure of \(S\) in the chosen separable closure of its fraction field, and let \(\bar i\) and \(\bar j\) be the inclusions of the special fibre and of the geometric generic fibre of this pullback.  Then, for \(L\) on \(\mathcal Y_\eta\) with pullback \(L_{\bar\eta}\) to the geometric generic fibre, the geometric nearby cycles are

\[
R\Psi_{\mathcal Y}(L)=\bar i^{*}R\bar j_*L_{\bar\eta}.
\]

This defines a functor \(R\Psi_{\mathcal Y}:D^b_c(\mathcal Y_{\eta,\acute et},A)\to
D(\mathcal Y_{s,\acute et},A)\). We suppress the inertia action, except where invariants under the wild inertia group of a finite extension of traits are taken below.

% object: example-nearby-cycles-node (blueprint/chap-Specialization-and-contracting-correspondences/sec-Specialization-to-the-normal-cone/example-nearby-cycles-node/_section.md)
\phantomsection\label{bp5}%
\begin{bpstatement}{Example~E.1}{ (A node)}
Nearby cycles see more than the special fibre when the family degenerates. For example, for \(\mathcal Y=V(xy-t)\subset\mathbf A^2_S\) the special fibre is the pair of axes \(xy=0\), while the generic fibre is smooth. To see what to expect, note that over \(\mathbf C\) the part of the fibre \(xy=\varepsilon\) near the origin is an annulus \(\{(x,\varepsilon/x)\}\), which retracts onto a circle that vanishes as \(\varepsilon\to0\). Correspondingly, in our setting the stalk of \(R\Psi_{\mathcal Y}(A)\) at the node is \(A\) in degree zero and \(A(-1)\) in degree one; the twist is that of \(H^1(\mathbf G_m)\) (see [\hyperref[bp100]{SGA7II}, Exposé XV]). At every other point of the special fibre, by contrast, the family is smooth over \(S\), and there the nearby cycles are \(A\) in degree zero, by smooth base change and the first calculation below.
\end{bpstatement}

% text: prose-E-03
\phantomsection\label{bp6}%
For a morphism \(g:\mathcal Y'\to\mathcal Y\) of separated \(S\)-schemes of finite type, there is a base change map for pullback

\[
g_s^*R\Psi_{\mathcal Y}
\longrightarrow R\Psi_{\mathcal Y'}g_\eta^*,
\]

which is invertible for smooth \(g\). For proper \(g\) there is a base change map for proper direct image

\[
Rg_{s*}R\Psi_{\mathcal Y'}
\xrightarrow{\sim}R\Psi_{\mathcal Y}Rg_{\eta*}.
\]

These base change maps commute with composition and with the specialization map defined below. Their invertibility follows from two base change theorems for torsion coefficients. First, proper base change [Stacks, \href{https://stacks.math.columbia.edu/tag/095T}{095T}] says that restriction to a fibre commutes with proper direct image. Second, smooth base change [Stacks, \href{https://stacks.math.columbia.edu/tag/0F09}{0F09}] (for sheaves [Stacks, \href{https://stacks.math.columbia.edu/tag/0EYU}{0EYU}]) says that pullback along a smooth map commutes with direct image along a qcqs morphism, when the coefficients are torsion of order invertible on the base. Suppose first that \(g\) is proper. The two direct images on the generic fibres agree because direct image is compatible with composition, and \(\bar i^*\) commutes with the proper direct image by proper base change. Suppose next that \(g\) is smooth. Since \(R\bar j_*\) is direct image along a qcqs morphism, we may apply smooth base change to it and then restrict to the special fibre. Smooth base change applies because the coefficients are annihilated by a power of \(\ell\), which is invertible on \(S\). The same base change maps appear in [\hyperref[bp100]{SGA7II}], Exposé XIII, (2.1.7.1)–(2.1.7.3).

We also use two calculations for a finite \(A\)-module \(M\). The first is \(R\Psi_S(M_{\eta})=M\). It holds because the integral closure \(\bar S\) of \(S\) in the separable closure is strictly henselian, so the stalk of \(R\bar j_*M\) at its closed point is \(R\Gamma(\bar\eta,M)=M\). The second is the usual calculation for the affine line with coefficients prime to the characteristic: \(H^r(\mathbf A^1_k,M)=0\) for \(r>0\) and \(H^0(\mathbf A^1_k,M)=M\). It follows from [Stacks, \href{https://stacks.math.columbia.edu/tag/03RR}{03RR}] applied to the compactification \(\mathbf A^1\subset\mathbf P^1\), which has genus zero and one point at infinity. To apply that result, note that every finite \(M\) is a finite direct sum of cyclic modules, and over the algebraically closed field \(k\) the cyclic constant coefficients are isomorphic to the corresponding roots of unity. By truncation the two calculations extend to a bounded complex \(M^\bullet\) of finite \(A\)-modules, viewed as a constant complex:

\[
R\Psi_S(M^\bullet_\eta)\simeq M^\bullet,
\qquad R\Gamma(\mathbf A^1_k,M^\bullet_{\mathbf A^1_k})\simeq M^\bullet.
\]

Finally, all complexes stay bounded and constructible. By the finiteness theorem for nearby cycles with torsion coefficients ([\hyperref[bp101]{SGA4½}], \emph{Théorèmes de finitude}, §3; [illusie94, §4.1, pp. 44–45]), \(R\Psi_{\mathcal Y}\) takes bounded constructible complexes of finite Tor dimension to complexes of the same kind. Finite Tor dimension is what we need when we take traces. The same holds for an arbitrary bounded constructible complex, without the Tor condition. To see this, we may forget the \(A\)-action, since the underlying torsion sheaves of \(A\)-linear nearby cycles are the usual nearby cycles. Bounded constructible complexes are stable under extensions and truncation, so the filtration by powers of \(\ell\) reduces us to constructible \(\mathbf F_\ell\)-sheaves. These have finite Tor dimension over \(\mathbf F_\ell\), and the finiteness theorem applies.

For a separated \(k\)-scheme \(Y\) of finite type, the constant family \(Y\times S\) has special fibre \(Y\). For \(L\) on \(Y\), let \(L_\eta\) be its pullback to the generic fibre.  Restricting the adjunction unit of \(\bar j\) on the pullback of \(L\) to the special fibre, we get the specialization map \(\sigma_Y(L)\) below.

% object: lemma-nearby-cycles-constant-family (blueprint/chap-Specialization-and-contracting-correspondences/sec-Specialization-to-the-normal-cone/lemma-nearby-cycles-constant-family/_section.md)
\phantomsection\label{bp7}%
\begin{bpstatementit}{Lemma~E.2}{ (Nearby cycles of a constant family)}
For a separated \(k\)-scheme \(Y\) of finite type and \(L\in\Dbc(Y,A)\), the specialization map

\[
\sigma_Y(L):L\longrightarrow
R\Psi_{Y\times S}(L_\eta)
\]

is an isomorphism. It is natural in \(L\), and it is compatible with the base change maps above in the following sense.

\begin{enumerate}
\item If \(f:Y'\to Y\) is proper and \(L'\in\Dbc(Y',A)\), then the composite of \(Rf_*\sigma_{Y'}(L')\) with the base change map for proper direct image is \(\sigma_Y(Rf_*L')\), where \((Rf_*L')_\eta\) is identified with \(Rf_{\eta*}(L'_\eta)\) by proper base change.
\item If \(f:Y'\to Y\) is smooth, then the composite of \(f^*\sigma_Y(L)\) with the base change map for pullback is \(\sigma_{Y'}(f^*L)\).
\end{enumerate}
\end{bpstatementit}

% proof: proof/lemma-nearby-cycles-constant-family
\begin{proof}[{Proof}]\phantomsection\label{bp8}%
The compatibilities (1) and (2) are formal. For the isomorphism we show that the cone \(C_Y(L)\) of \(\sigma_Y(L)\) vanishes. To do this we use that, like nearby cycles, the cone commutes with smooth pullback and proper direct image. Smooth pullback lets us check vanishing locally, and it also shows that the cone vanishes generically, by reduction to constant coefficients and the calculation \(R\Psi_S(M_\eta)=M\). Proper direct image along a compactified linear projection then shows that the remaining support is empty, by induction on the dimension of an affine space containing \(Y\).

By the comparison for constructible complexes, we may work on the small étale site. There \(\sigma\) is \(\bar i^*\) of the unit of \(\bar j^*\dashv R\bar j_*\) and the base change maps are \(\bar i^*\) of mates, so (1), (2) and naturality follow from the triangle identities.

We take cofibres in the stable \(\infty\)-category underlying the derived category, where they are functorial, unlike cones in a triangulated category, and set \(C_Y(L)=\operatorname{cofib}(\sigma_Y(L))\).  This is an exact functor with bounded constructible values.  By (1) and (2), the base change maps yield

\[
C_Y(Rf_*L')\simeq Rf_*C_{Y'}(L')\quad(f\text{ proper}),\qquad
C_{Y'}(f^*L)\simeq f^*C_Y(L)\quad(f\text{ smooth}).
\]

By the base change isomorphism for étale maps, the question is local on \(Y\).  For a closed immersion \(Y\hookrightarrow\mathbf A^n\), the base change isomorphism for proper maps identifies the direct image of \(C_Y(L)\) with \(C_{\mathbf A^n}\) of the direct image of \(L\), and since direct image along a closed immersion is conservative, it suffices to treat affine spaces.

For \(Y=\mathbf A^n\), the complex \(C_Y(L)\) vanishes on a dense open subset.  Indeed, by the truncation triangles we may assume that \(L\) is a sheaf. After shrinking, it is locally constant, and after pulling back to a finite étale cover, it is constant. The pullback does no harm: \(C\) commutes with étale pullback, so vanishing can be checked after a surjective étale pullback.  For a constant sheaf \(M\) on a smooth \(k\)-scheme, the base change isomorphism for its smooth structure map identifies \(C\) with the pullback of \(C_{\operatorname{Spec}k}(M)\), which is zero because \(R\Psi_S(M_\eta)=M\).

We argue by induction on \(n\), the case of a point being the calculation above.  By the preceding paragraph, the closure \(T\) of the support of \(C_{\mathbf A^n}(L)\) is a proper closed subset.  By the linear form of Noether normalization there is a linear projection \(q:\mathbf A^n\to\mathbf A^{n-1}\) that is finite on \(T\): we choose its kernel line so that it avoids the intersection of the projective closure of \(T\) with the hyperplane at infinity.  Compactifying that line, we get

\[
\bar q:\bar Y=\mathbf P^1\times\mathbf A^{n-1}
\longrightarrow\mathbf A^{n-1},
\qquad v:\mathbf A^n\hookrightarrow\bar Y,
\qquad \bar L=v_!L.
\]

We compactified the line so that the support becomes finite over \(\mathbf A^{n-1}\) and can be detected by direct image there. As \(C\) commutes with restriction to the open subscheme \(\mathbf A^n\) of \(\bar Y\), the complex \(C_{\bar Y}(\bar L)\) is supported on \(T\cup(\bar Y\setminus\mathbf A^n)\).  This union is closed in \(\bar Y\) and finite over \(\mathbf A^{n-1}\). Indeed, \(T\) is finite, hence proper, over \(\mathbf A^{n-1}\), and therefore closed in \(\bar Y\); the boundary \(\bar Y\setminus\mathbf A^n\) is a section.  By the base change isomorphism for proper maps and the induction hypothesis,

\[
R\bar q_*C_{\bar Y}(\bar L)
\simeq C_{\mathbf A^{n-1}}(R\bar q_*\bar L)=0.
\]

Direct image is conservative on complexes supported on a closed subscheme finite over the base, since on a geometric stalk the direct image of each cohomology sheaf is the finite product of its stalks on the finite fibre.  Therefore \(C_{\bar Y}(\bar L)=0\), and restricting to \(\mathbf A^n\) gives \(C_{\mathbf A^n}(L)=0\).
\end{proof}

% text: prose-E-04
\phantomsection\label{bp9}%
The contracting argument needs a complex on the normal cone \(N_ZX\) that still remembers \(K\) near \(Z\). Pulling \(K|_Z\) back along the projection \(N_ZX\to Z\) is not enough, since it forgets the directions leaving \(Z\). For example, take \(X=\mathbf A^1\), \(Z\) the origin, and \(K=j_!A\) for \(j:\mathbf A^1\setminus\{0\}\hookrightarrow\mathbf A^1\). Then \(K|_Z=0\), so the pullback is zero. By contrast, the deformation to the normal cone is \(\operatorname{Spec}k[x,t,x/t]\) over the trait, that is, the constant family \(\mathbf A^1\times S\) in the coordinate \(x'=x/t\), and on its generic fibre \(K\) is \(j_!A\) for \(x'\neq0\). Hence, by Lemma \hyperref[bp7]{E.2}, the nearby cycles are \(j_!A\) on \(N_ZX=\mathbf A^1\), which is not zero. This construction is the specialization functor \(\operatorname{Sp}_{X,Z}\) defined below, and by Verdier's theorem (Theorem \hyperref[bp24]{E.8}) it still recovers \(K|_Z\) on the zero section. Theorem \hyperref[bp91]{E.30} uses both properties.

Let \(X\) be a separated \(k\)-scheme of finite type, and let \(i:Z\hookrightarrow X\) be a closed subscheme with ideal \(I\). The deformation to its normal cone, over the chosen trait, is

\[
\mathcal D_ZX=
\operatorname{Spec}_{X\times S}
\bigl(\mathcal O_{X\times S}[I/t]\bigr).
\]

Here \(\mathcal O_{X\times S}[I/t]\) is the base change to \(S\) of the subalgebra of \(\mathcal O_X[t,t^{-1}]\) generated by \(\mathcal O_X[t]\) and the fractions \(a/t\) with \(a\in I\).  Its generic fibre is \(X_\eta\) and its special fibre is

\[
N_ZX=\operatorname{Spec}_Z
\left(\bigoplus_{r\geq0}I^r/I^{r+1}\right).
\]

The quotient by \(I\) and all \(a/t\) defines a closed immersion \(\widetilde z:Z\times S\hookrightarrow\mathcal D_ZX\), whose special fibre \(z:Z\hookrightarrow N_ZX\) is the zero section.  For \(L\) on \(X\), let \(L_\eta\) be its pullback to the generic fibre \(X_\eta\) of \(\mathcal D_ZX\), and define the specialization functor \(\operatorname{Sp}_{X,Z}:D^b_c(X_{\acute et},A)\to
D^b_c((N_ZX)_{\acute et},A)\) by

\[
\operatorname{Sp}_{X,Z}(L)=R\Psi_{\mathcal D_ZX}(L_\eta).
\]

The base change map for pullback along \(\widetilde z\), followed by the inverse of the isomorphism of Lemma \hyperref[bp7]{E.2}, defines a natural transformation

\[
\tau_{X,Z}(L):z^*\operatorname{Sp}_{X,Z}(L)
\longrightarrow R\Psi_{Z\times S}((i^*L)_\eta)
\xrightarrow{\sim}i^*L.
\tag{E.1}
\]

We use this map, in this direction, to specialize trace classes.

% object: example-normal-cone-point (blueprint/chap-Specialization-and-contracting-correspondences/sec-Specialization-to-the-normal-cone/example-normal-cone-point/_section.md)
\phantomsection\label{bp10}%
\begin{bpstatement}{Example~E.3}{ (Normal cones of a point)}
If \(X\) is smooth at a point \(Z=\{y\}\), then \(\bigoplus_rI^r/I^{r+1}\) is the symmetric algebra of \(I/I^2\) and \(N_ZX\) is the tangent space at \(y\). At a singular point, however, the normal cone is the tangent cone, and it need not be reduced. For example, for the cusp \(X=V(y^2-x^3)\subset\mathbf A^2\) and \(Z\) the origin, with \(I=(x,y)\), the graded algebra is \(k[\xi,\eta]/(\eta^2)\), where \(\xi,\eta\) are the classes of \(x,y\) in \(I/I^2\) and \(\eta^2\) is the lowest-degree part of \(y^2-x^3\); so \(N_ZX\) is a double line. Similarly, for a nonreduced point \(Z=V(x^m)\subset\mathbf A^1\), each \(I^r/I^{r+1}=(x^{mr})/(x^{m(r+1)})\) is free of rank one over \(\mathcal O_Z=k[x]/(x^m)\), and \(N_ZX=Z\times\mathbf A^1\) is a line thickened along \(Z\); this is the case treated in Lemma \hyperref[bp18]{E.6}.
\end{bpstatement}

% text: prose-E-05
\phantomsection\label{bp11}%
We will apply the base change maps for nearby cycles to the deformations of different pairs, and for this we need to know how the deformation depends on the pair.

% object: lemma-functoriality-deformation (blueprint/chap-Specialization-and-contracting-correspondences/sec-Specialization-to-the-normal-cone/lemma-functoriality-deformation/_section.md)
\phantomsection\label{bp12}%
\begin{bpstatementit}{Lemma~E.4}{ (Functoriality of the deformation)}
Let \(f:X'\to X\) be a morphism of separated \(k\)-schemes of finite type, let \(Z'\subset X'\) be a closed subscheme with ideal \(I'\), and suppose that \(f^{-1}(I)\mathcal O_{X'}\subset I'\).  Then \(f\) induces \(\widetilde f:\mathcal D_{Z'}X'\to\mathcal D_ZX\). If \(Z'=f^{-1}Z\) schematically, then

\[
\mathcal D_{Z'}X'\longrightarrow X'\times_X\mathcal D_ZX
\]

is a closed immersion, and it is an isomorphism if \(f\) is flat. Moreover, \(\widetilde f_s^{-1}(Z)=Z'\) schematically. In this case \(\widetilde f\) is proper when \(f\) is proper and smooth when \(f\) is smooth.
\end{bpstatementit}

% proof: proof/lemma-functoriality-deformation
\begin{proof}[{Proof}]\phantomsection\label{bp13}%
All assertions are statements about the algebras \(\mathcal O[I/t]\), and we check them on generators. The first morphism is defined by \(a/t\mapsto f^\#(a)/t\), which makes sense because \(f^\#(a)\in I'\) for \(a\in I\). If \(I'=I\mathcal O_{X'}\), the natural map from the base change of the Rees algebra onto \(\mathcal O_{X'\times S}[I'/t]\) is surjective, and it is injective if \(f\) is flat.  On the associated graded algebras, the generators of \(\operatorname{gr}_I\) of positive degree generate the ideal of elements of positive degree of \(\operatorname{gr}_{I'}\).  The quotient by them is \(\mathcal O_{Z'}\), and this proves the assertion about the zero section.  If \(f\) is proper, then \(\widetilde f\) is the closed immersion followed by the base change of \(f\), which is proper. If \(f\) is smooth, it is flat, so \(\widetilde f\) is the base change \(X'\times_X\mathcal D_ZX\to\mathcal D_ZX\) of \(f\), which is smooth.
\end{proof}

% text: prose-E-06
\phantomsection\label{bp14}%
Theorem \hyperref[bp24]{E.8} asserts that \(\tau_{X,Z}\) is an isomorphism, that is, that the cofibre

\[
C_{X,Z}(L)=\operatorname{cofib}\bigl(
\tau_{X,Z}(L)\bigr)
\]

vanishes. As for \(C_Y\), we compute the cofibre in the stable \(\infty\)-category, so that \(C_{X,Z}\) is an exact functor. The proof has three steps.  First, the base change maps of the next lemma reduce us to the case \(i^*L=0\) and, after a blow-up, to the case of an effective Cartier divisor.  Then, for a principal divisor in a normal scheme, an explicit contraction of the affine line shows that \(C_{X,Z}(L)\) vanishes at the generic points of \(Z\).  Finally, after an embedding into affine space, projection along a line and induction on the dimension show that the remaining support is empty.

% object: lemma-base-change-cone-x-z (blueprint/chap-Specialization-and-contracting-correspondences/sec-Specialization-to-the-normal-cone/lemma-base-change-cone-x-z/_section.md)
\phantomsection\label{bp15}%
\begin{bpstatementit}{Lemma~E.5}{ (Base change for the cone of $\tau_{X,Z}$)}
Let \(f:X'\to X\) be a morphism of separated \(k\)-schemes of finite type, let \(Z'=f^{-1}Z\) schematically, and let \(f_Z:Z'\to Z\) be the restriction of \(f\).  There are isomorphisms

\[
\begin{aligned}
C_{X,Z}(Rf_*L')&\simeq Rf_{Z*}C_{X',Z'}(L')
&&(f\text{ proper}),\\
C_{X',Z'}(f^*L)&\simeq f_Z^*C_{X,Z}(L)
&&(f\text{ smooth}).
\end{aligned}
\]

Moreover, \(C_{X,Z}(i_*M)=0\) for every \(M\in D^b_c(Z_{\acute et},A)\). Consequently, writing \(j:X\setminus Z\hookrightarrow X\), we have

\[
C_{X,Z}(j_!j^*L)\xrightarrow{\sim}C_{X,Z}(L).
\]
\end{bpstatementit}

% proof: proof/lemma-base-change-cone-x-z
\begin{proof}[{Proof}]\phantomsection\label{bp16}%
Both the source and the target of \(\tau\) are compatible with proper direct image and smooth pullback, so the cone is as well. To prove the vanishing on \(i_*M\), we reduce to the pair \((Z,Z)\), where the deformation is a constant family.

By Lemma \hyperref[bp12]{E.4}, \(f\) induces a morphism \(\widetilde f\) of deformations, which is proper if \(f\) is proper and smooth if \(f\) is smooth. To this morphism we apply the base change maps of nearby cycles, together with proper or smooth base change for the restriction to the zero section.  The resulting squares for \(\tau\) then commute because the base change map for pullback is compatible with composition, and their cofibres give the two isomorphisms. In the special case \(f=i\), the source pair is \((Z,Z)\) and \(\mathcal D_ZZ=Z\times S\), so \(\tau_{Z,Z}\) is the inverse of the isomorphism of Lemma \hyperref[bp7]{E.2}. Hence \(C_{Z,Z}=0\), and since \(i\) is proper, the first isomorphism gives \(C_{X,Z}(i_*M)=0\). The last assertion then follows by applying the exact functor \(C_{X,Z}\) to \(j_!j^*L\to L\to i_*i^*L\).
\end{proof}

% text: prose-E-07
\phantomsection\label{bp17}%
To prove the vanishing at the generic points of \(Z\), we first treat a nonreduced point on the affine line, where an explicit contraction is available.

% object: lemma-case-z-v-x-m (blueprint/chap-Specialization-and-contracting-correspondences/sec-Specialization-to-the-normal-cone/lemma-case-z-v-x-m/_section.md)
\phantomsection\label{bp18}%
\begin{bpstatementit}{Lemma~E.6}{ (The case $Z=V(x^m)\subset\mathbf A^1$)}
For \(X=\mathbf A^1_k\), \(Z=V(x^m)\) with \(m\geq1\), and every finite \(A\)-module \(M\), one has \(C_{X,Z}(M_X)=0\).
\end{bpstatementit}

% proof: proof/lemma-case-z-v-x-m
\begin{proof}[{Proof}]\phantomsection\label{bp19}%
We give a homotopy argument. The contraction \(x\mapsto ax\) joins the identity (\(a=1\)) to a map that factors through \(Z\) (\(a=0\)), and the endomorphisms of \(z^*\operatorname{Sp}_{X,Z}(M_X)\) that these maps induce form a family over \(\mathbf A^1\). Since \(\mathbf A^1\) has the cohomology of a point, the family is constant, and comparing its two ends shows that \(\tau_{X,Z}(M_X)\) has an inverse.

Consider the contraction \(h:X\times\mathbf A^1\to X\), \(h(x,a)=ax\), and write \(h_a=h(-,a):X\to X\) for \(a\in k\).  It maps \(Z\times\mathbf A^1\) schematically into \(Z\), since \(h^*(x^m)=a^m x^m\in(x^m)\), and by Lemma \hyperref[bp12]{E.4}, whose hypothesis is this inclusion of ideals, all \(h_a\) lift to the deformations.  Moreover, \(h_1=\operatorname{id}\), and \(h_0\) factors through \(Z\) and restricts to the identity on \(Z_{\mathrm{red}}\).  Set \(N=z^*\operatorname{Sp}_{X,Z}(M_X)\). The base change maps for pullback along these lifts induce endomorphisms \(\phi_a:N\to N\), with \(\phi_1=\operatorname{id}_N\).

There is one endomorphism \(\phi\) of the pullback of \(N\) to \(Z\times\mathbf A^1\) whose fibre at \(a\) is \(\phi_a\).  To construct it, we apply the base change map for pullback along \(h:X\times\mathbf A^1\to X\) and then invert the base change map for the projection to \(X\), which is smooth; by flat base change, the deformation of \(X\times\mathbf A^1\) is \(\mathcal D_ZX\times\mathbf A^1\).  The restrictions of \(h\) and of the projection to \(Z\times\mathbf A^1\) induce the same map to \(Z\) on reduced schemes, so invariance under nilpotent thickenings identifies \(h_Z^*N\) with \(\operatorname{pr}_Z^*N\), naturally in \(N\), and we may view the family as an endomorphism of a single pullback.  By compatibility with composition, its fibres are the endomorphisms \(\phi_a\) constructed above.

The family \(\phi\) is constant; this is the algebraic form of homotopy invariance. Since \(Z_{\mathrm{red}}\) is a point, invariance under nilpotent thickenings identifies \(D^b_c(Z_{\acute et},A)\) with the bounded derived category of finite \(A\)-modules.  Pullback along the projection \(p:Z\times\mathbf A^1\to Z\) is fully faithful on this category. Indeed, the calculation for the affine line and the hypercohomology spectral sequence show that \(R\Gamma(\mathbf A^1,p^*B)=B\) for every such bounded complex \(B\), and full faithfulness follows by adjunction. Thus \(\phi\) is the pullback of a single endomorphism of \(N\). Its two restrictions \(\phi_0\) and \(\phi_1\) therefore coincide, and \(\phi_0=\operatorname{id}_N\).

Because \(h_0\) factors through \(Z\), we can also express \(\phi_0\) in terms of \(\tau_{X,Z}(M_X)\).  Let \(\alpha:M_Z\to N\) be the restriction of the unit of nearby cycles for the constant sheaf on \(\mathcal D_ZX\).  By naturality, and since \(h_0\) factors as \(X\to Z\hookrightarrow X\), we get

\[
\phi_0=\alpha\tau_{X,Z}(M_X),\qquad
\tau_{X,Z}(M_X)\alpha=\operatorname{id}_{M_Z}.
\]

The second equality holds because the composite is the specialization map for the constant family \(Z\times S\), followed by its fixed inverse.  Since \(\phi_0=\phi_1=\operatorname{id}_N\), the maps \(\alpha\) and \(\tau_{X,Z}(M_X)\) are inverse isomorphisms.
\end{proof}

% text: prose-E-08
\phantomsection\label{bp20}%
Near a generic point of a principal divisor in a normal scheme, the pair \((X,Z)\) is a smooth pullback of \((\mathbf A^1,V(x^m))\), so by Lemma \hyperref[bp18]{E.6} \(C_{X,Z}\) vanishes there on constant coefficients. The next lemma extends this to all sheaves that vanish on \(Z\).

% object: lemma-vanishing-generic-points-principal-divisor (blueprint/chap-Specialization-and-contracting-correspondences/sec-Specialization-to-the-normal-cone/lemma-vanishing-generic-points-principal-divisor/_section.md)
\phantomsection\label{bp21}%
\begin{bpstatementit}{Lemma~E.7}{ (Vanishing at the generic points of a principal divisor)}
Suppose that \(X\) is normal, that \(Z\subset X\) is an effective principal Cartier divisor, and that \(i^*L=0\).  Then \(C_{X,Z}(L)\) vanishes at every generic point of \(Z\).
\end{bpstatementit}

% proof: proof/lemma-vanishing-generic-points-principal-divisor
\begin{proof}[{Proof}]\phantomsection\label{bp22}%
For a constant sheaf extended by zero, a smooth map \(X\to\mathbf A^1\) near the generic point reduces the claim to Lemma \hyperref[bp18]{E.6}. Near the generic point, a general sheaf is \(j_!G\) with \(G\) locally constant. After a finite étale cover, \(G\) becomes constant, while the cokernel of the unit of the cover is again of the form \(j_!G'\). Passing to this cokernel shifts the cohomological degree by one, so an induction on the degree proves the claim for all \(G\) at once.

Since the question is local at a generic point \(\zeta\) of \(Z\), we may discard the other irreducible components of \(X\) and \(Z\).  Consider first \(L=j_!M_U\), where \(U=X\setminus Z\) and \(M\) is a finite \(A\)-module. Since \(X\) is normal, it is regular in codimension one.  After shrinking around \(\zeta\), both \(X\) and \(Z_{\mathrm{red}}\) are smooth over \(k\), and a local equation \(h\) of \(Z_{\mathrm{red}}\) defines a smooth map \(h:X\to\mathbf A^1\).  The local ring \(\mathcal O_{X,\zeta}\) is a discrete valuation ring with uniformizer \(h\), so a local equation \(g\) of \(Z\) can be written \(g=wh^m\) with \(m\geq1\) and \(w\) a unit of \(\mathcal O_{X,\zeta}\).  The function \(w\) is regular and invertible on a neighbourhood of \(\zeta\), and there the ideal of \(Z\) is \((h^m)\), that is, \(Z=h^{-1}(V(x^m))\) schematically.  By the base change isomorphism for smooth maps and Lemma \hyperref[bp18]{E.6}, we see that \(C_{X,Z}(M_X)=0\) there, and the last assertion of Lemma \hyperref[bp15]{E.5} yields the same vanishing for \(j_!M_U\).

For general \(L\), truncation triangles reduce us to a constructible sheaf \(F\) with \(F|_Z=0\).  Shrinking around \(\zeta\), we may assume that \(F=j_!G\) for a finite locally constant sheaf \(G\) on \(U\), since the closure in \(X\) of the locus in \(U\) where \(F\) is not locally constant cannot contain a generic point of the divisor \(Z\). We prove by induction on \(r\), simultaneously for all such normal pairs and all \(G\), that \(H^r(C_{X,Z}(j_!G))\) vanishes generically. As nearby cycles are left \(t\)-exact and \(i^*j_!G=0\), we have

\[
C_{X,Z}(j_!G)=
z^*\operatorname{Sp}_{X,Z}(j_!G)[1]\in D^{\geq-1}.
\]

Therefore the statement holds for every \(r<-1\), which starts the induction.

For the induction step, we compare \(F=j_!G\) with a sheaf to which the first part of the proof applies. To construct it, choose a finite étale surjection \(f_U:U'\to U\) trivializing \(G\), and let \(f:X'\to X\) be the normalization of \(X\) in \(U'\); it is finite since \(X\) is excellent, and \(X'\) is normal.  With \(Z'=f^{-1}Z\) and \(j':U'\hookrightarrow X'\), the sheaf \(f^*F=j'_!f_U^*G\) is \(j'_!\) of a sheaf that is constant on each component of \(U'\), and \(Z'\) is a principal divisor in the normal scheme \(X'\). Since \(f\) is finite, the points of \(Z'\) over a generic point of \(Z\) are generic points of \(Z'\), so the first part of the proof, applied to \((X',Z')\), and the base change isomorphism for proper maps show that \(C_{X,Z}(f_*f^*F)\) vanishes generically on \(Z\).

On a geometric stalk, the unit of the finite étale cover is the diagonal inclusion into a nonempty finite product, so it is injective and there is a short exact sequence

\[
0\longrightarrow F\longrightarrow f_*f^*F
\longrightarrow j_!G'\longrightarrow0,
\]

where \(G'\) is again finite locally constant on \(U\).  Applying \(C_{X,Z}\) and taking the long exact sequence at a generic point where the middle term vanishes, we find

\[
H^r(C_{X,Z}(F))\simeq
H^{r-1}(C_{X,Z}(j_!G')).
\]

By the induction hypothesis, the right side vanishes.
\end{proof}

% text: prose-E-09
\phantomsection\label{bp23}%
On \(D^b_{\mathrm{ctf}}(X,A)\) we define the functor \(\operatorname{Sp}_{X,Z}\) and the map (E.1) by transport of the étale constructions above along the comparison for constructible complexes with finite coefficients (Theorem \bpout{2.43}).

% object: theorem-verdier-s-theorem-specialization-zero (blueprint/chap-Specialization-and-contracting-correspondences/sec-Specialization-to-the-normal-cone/theorem-verdier-s-theorem-specialization-zero/_section.md)
\phantomsection\label{bp24}%
\begin{bpstatementit}{Theorem~E.8}{ (Verdier's theorem on specialization to the zero section)}
For every separated \(k\)-scheme \(X\) of finite type, every closed subscheme \(Z\subset X\), and \(K\in D^b_{\mathrm{ctf}}(X,A)\), the map (E.1)

\[
\tau_{X,Z}(K):z^*\operatorname{Sp}_{X,Z}(K)\longrightarrow i^*K
\]

is an isomorphism.  In particular, \(i^*K=0\) implies \(z^*\operatorname{Sp}_{X,Z}(K)=0\).
\end{bpstatementit}

% proof: proof/theorem-verdier-s-theorem-specialization-zero
\begin{proof}[{Proof}]\phantomsection\label{bp25}%
The main step is that \(C_{X,Z}(L)=0\) for every \(L\in
D^b_c(X_{\acute et},A)\), which the last paragraph transports to \(D^b_{\mathrm{ctf}}\).  We need the vanishing for all bounded constructible complexes, because the cokernels \(j_!G'\) in the dimension shifting of Lemma \hyperref[bp21]{E.7} may have infinite Tor dimension.

By Lemma \hyperref[bp15]{E.5}, we may replace \(L\) by \(j_!j^*L\) and so assume that \(i^*L=0\).  We may then replace \(X\) by its blow-up \(\pi:X'=\operatorname{Bl}_ZX\to X\) along \(Z\). The reason is that a complex that vanishes on the centre is recovered from its pullback to the blow-up, that is, the unit \(L\to R\pi_*\pi^*L\) is an isomorphism. Indeed, away from \(Z\) this holds because \(\pi\) is an isomorphism there, and on \(Z\) it holds because proper base change identifies the target with the cohomology of a fibre with zero coefficients.  Using this, the base change isomorphism for proper maps gives \(C_{X,Z}(L)\simeq C_{X,Z}(R\pi_*\pi^*L)\simeq R\pi_{Z*}C_{X',\pi^{-1}Z}(\pi^*L)\), and it suffices to treat the pair \((X',\pi^{-1}Z)\).  Since the pullback of the ideal is invertible on the blow-up, we are reduced to the case of a Cartier divisor.

To have linear projections available, we move the pair into affine space. By the base change isomorphism for open immersions the assertion is local, so we may assume that \(X\) is affine and \(Z=V(h)\) for a function \(h\).  Choosing algebra generators for \(\mathcal O(X)\) and appending \(h\) as the last coordinate gives a closed immersion

\[
\iota:X\hookrightarrow\mathbf A^n,\qquad
Z=\iota^{-1}(\mathbf A^{n-1}\times\{0\})
\]

schematically.  The base change isomorphism for proper maps identifies \(R\iota_{Z*}C_{X,Z}(L)\) with \(C_{\mathbf A^n,\mathbf A^{n-1}\times\{0\}}(R\iota_*L)\), and since direct image along a closed immersion is conservative, it suffices to treat the pair \((\mathbf A^n,\mathbf A^{n-1}\times\{0\})\). This is a principal divisor in a normal scheme, so Lemma \hyperref[bp21]{E.7} applies to it.

We argue by induction on \(n\geq1\), again assuming that \(L|_Z=0\). By Lemma \hyperref[bp21]{E.7}, the closure \(Y\) of the support of \(C_{X,Z}(L)\) is a proper closed subset of the hyperplane \(Z\); for \(n=1\) it is empty, and the assertion follows.  For \(n>1\), choose a line \(\lambda\) in the vector space \(Z\) such that the linear projection

\[
q:X\longrightarrow X'=X/\lambda\simeq\mathbf A^{n-1}
\]

is finite on \(Y\).  Such a line exists: if the projective point \(\mathbf P(\lambda)\) lies outside the part at infinity of the projective closure of \(Y\) in \(\mathbf P(Z\oplus k)\), the usual proof of linear Noether normalization shows that the projection is finite on \(Y\). Set \(Z'=Z/\lambda\), a coordinate hyperplane in \(X'\), and compactify the \(\lambda\)-direction:

\[
\bar q:\bar X=\mathbf P^1\times X'\longrightarrow X',\qquad
\bar Z=\bar q^{-1}(Z'),\qquad
v:X\hookrightarrow\bar X,\qquad \bar L=v_!L.
\]

Let \(\bar q_Z:\bar Z\to Z'\) be the restriction of \(\bar q\). By the base change isomorphism for proper maps and the induction hypothesis,

\[
R\bar q_{Z*}C_{\bar X,\bar Z}(\bar L)
\simeq C_{X',Z'}(R\bar q_*\bar L)=0.
\]

Since the cone commutes with restriction to the open subscheme \(X\) of \(\bar X\), the complex \(C_{\bar X,\bar Z}(\bar L)\) is supported on \(Y\cup(\bar Z\setminus Z)\).  This is a closed subset finite over \(Z'\). Indeed, \(Y\) is finite over \(Z'\) and therefore remains closed in \(\bar Z\), and \(\bar Z\setminus Z\) is the section at infinity of \(\bar Z\to Z'\). As in Lemma \hyperref[bp7]{E.2}, direct image along a finite map is conservative on complexes with this support. Hence \(C_{\bar X,\bar Z}(\bar L)=0\), and restriction to \(X\) gives \(C_{X,Z}(L)=0\) on the étale site.

By the comparison for finite coefficients, \(K\) and its maps descend to the classical étale constructible category, where \(C_{X,Z}\) vanishes.  The comparison is exact and fully faithful, and it identifies pullback, proper direct image and their adjunction units, so the étale map (E.1) corresponds to the map in the statement, which is therefore an isomorphism.
\end{proof}


% object: sec-Kunneth-formula-for-nearby-cycles (blueprint/chap-Specialization-and-contracting-correspondences/sec-Kunneth-formula-for-nearby-cycles/_section.md), content hidden
\subsection*{\texorpdfstring{E.2\quad Künneth formula for nearby cycles}{E.2 Künneth formula for nearby cycles}}
\phantomsection\label{bp26}%
\addcontentsline{toc}{subsection}{\texorpdfstring{E.2\quad Künneth formula for nearby cycles}{E.2 Künneth formula for nearby cycles}}
% text: prose-E-10
\phantomsection\label{bp27}%
To specialize a trace class we have to specialize the external product \(K\boxtimes \mathbb DK\), so we need the comparison map from the external product of nearby cycles to the nearby cycles of the external product to be an isomorphism. We prove this on the classical étale site by the dimension argument of [\hyperref[bp99]{Ill94}, Theorem 4.7, pp. 51–53].

In this section and the next three, \(S\) is an excellent strictly henselian trait, as announced at the start of this appendix, with separably closed residue field \(k\), generic point \(\eta\), and a chosen geometric generic point \(\bar\eta\).  Let \(p\) be the characteristic exponent of \(k\).  The coefficients lie in a finite commutative \(\ell\)-primary ring \(A\), with \(\ell\) invertible on \(S\).  On the geometric fibres we use the ordinary external Künneth formula, Proposition \bpout{8.18}. It applies over the separably closed field \(k\), since passing to an algebraic closure is a purely inseparable extension and does not change the étale topos.

We use two further facts about traits, both from [\hyperref[bp101]{SGA4½}], \emph{Théorèmes de finitude}, §§3.4 and 3.7. First, geometric nearby cycles are invariant under extension of traits, with the corresponding restriction of the inertia action (loc. cit., 3.7). Second, the calculation in loc. cit., 3.4, describes them by restriction to the generic point of a coordinate line, as follows. Let \(s'\) be a geometric generic point of the special fibre \(\mathbf A^1_k\) of \(\mathbf A^1_S\), and let \(S'=(\mathbf A^1_S)_{(s')}\) be the strict localization of \(\mathbf A^1_S\) at \(s'\).  This is a strictly henselian trait with the same uniformizer as \(S\).  With compatible geometric generic points, its inertia group fits into an exact sequence

\[
1\longrightarrow P\longrightarrow I'\longrightarrow I
\longrightarrow1,
\]

where \(P\) is pro-\(p\) (and is trivial when \(p=1\)). For \(a:W\to\mathbf A^1_S\), put \(W'=W\times_{\mathbf A^1_S}S'\), and let \(N'\) be the pullback of \(N\in D^b_{\mathrm{ctf}}(W_\eta,A)\).  There is a canonical identification

\[
(R\Psi_{W/S}N)|_{W_s\times_{\mathbf A^1_k}s'}
\simeq (R\Psi_{W'/S'}N')^P.
\tag{E.2}
\]

Here \((-)^P=R\Gamma(P,-)\) is exact on the discrete torsion coefficients prime to \(p\) under consideration.  The formula comes from factoring the map from the geometric generic fibre through the geometric generic point of \(S'\), which lies over \(\bar\eta\). The Galois group of the second step is \(P\). It is pro-\(p\) because the common uniformizer identifies the tame inertia quotients. The formula is compatible with the base change maps for pullback and with the product maps obtained from \(\bar i^*R\bar j_*\). To see this, we factor the maps of generic fibres in the same way and use that the adjunctions for direct image are compatible with composition.

The functor of \(P\)-invariants commutes with tensoring by an object on which \(P\) acts trivially: for a complex \(Q\) with continuous \(P\)-action and a complex \(B\) with trivial \(P\)-action, the projection map is an isomorphism

\[
(Q\mathbin{\otimes_A^{\mathbf L}}B)^P
\simeq Q^P\mathbin{\otimes_A^{\mathbf L}}B.
\tag{E.3}
\]

Indeed, \(P\)-invariants are exact on discrete torsion \(A\)-modules.  For a finite quotient \(G\) of \(P\), the invariants are the image of the averaging idempotent \(|G|^{-1}\sum_{g\in G}g\), since the \(p\)-power \(|G|\) is invertible in \(A\), and this idempotent commutes with tensoring by \(B\).  By continuity we pass from the finite quotients to \(P\), and resolving \(B\) by flat \(A\)-modules gives the derived identity.

For \(X,Y\) separated of finite type over \(S\), \(L\) on \(X_\eta\) and \(M\) on \(Y_\eta\), the comparison map \(\Pi_{L,M}\) below is the restriction to the special fibre of the map

\[
R\bar j_{X*}L_{\bar\eta}\boxtimes
R\bar j_{Y*}M_{\bar\eta}
\longrightarrow
R\bar j_{X\times Y,*}(L_{\bar\eta}\boxtimes M_{\bar\eta})
\]

defined by the base change maps for pullback and the lax monoidal structure of direct image.

% object: proposition-external-kunneth-formula-nearby-cycles (blueprint/chap-Specialization-and-contracting-correspondences/sec-Kunneth-formula-for-nearby-cycles/proposition-external-kunneth-formula-nearby-cycles/_section.md)
\phantomsection\label{bp28}%
\begin{bpstatementit}{Proposition~E.9}{ (External Künneth formula for nearby cycles)}
Let \(X,Y\) be separated of finite type over \(S\), and let \(L\in D^b_{\mathrm{ctf}}(X_\eta,A)\) and \(M\in D^b_{\mathrm{ctf}}(Y_\eta,A)\).  The comparison map

\[
\Pi_{L,M}:R\Psi_{X/S}L\boxtimes R\Psi_{Y/S}M
\longrightarrow
R\Psi_{(X\times_SY)/S}(L\boxtimes M)
\tag{E.4}
\]

is an isomorphism.
\end{bpstatementit}

% text: prose-E-11
\phantomsection\label{bp29}%
The proof has the same shape as that of Lemma \hyperref[bp7]{E.2}: the cone of \(\Pi_{L,M}\) has no global sections and finite support, so it vanishes. When \(X\) and \(Y\) are projective, the ordinary Künneth formula on the geometric generic fibres gives zero global sections, once we know how \(\Pi\) behaves under proper direct image. Restricting to the generic point of a coordinate line by (E.2) lowers the dimension of the generic fibre, and induction on the dimension shows that the support is finite. The next four lemmas supply these ingredients.

% object: lemma-kunneth-map-proper-direct-image (blueprint/chap-Specialization-and-contracting-correspondences/sec-Kunneth-formula-for-nearby-cycles/lemma-kunneth-map-proper-direct-image/_section.md)
\phantomsection\label{bp30}%
\begin{bpstatementit}{Lemma~E.10}{ (The Künneth map and proper direct image)}
Let \(f:X\to X_1\) and \(g:Y\to Y_1\) be proper morphisms of separated \(S\)-schemes of finite type, and let \(L\in D^b_{\mathrm{ctf}}(X_\eta,A)\) and \(M\in D^b_{\mathrm{ctf}}(Y_\eta,A)\). Identify direct images of external products by the ordinary Künneth formula, and direct images of nearby cycles by the base change isomorphism for proper maps. Then \(R(f_s\times g_s)_*\Pi_{L,M}\) becomes \(\Pi_{Rf_{\eta*}L,Rg_{\eta*}M}\).
\end{bpstatementit}

% proof: proof/lemma-kunneth-map-proper-direct-image
\begin{proof}[{Proof}]\phantomsection\label{bp31}%
Write the two product maps in terms of the tensor maps of \(R\bar j_*\), commute proper direct image past these maps on the generic fibres, and take the adjoints of the ordinary Künneth arrows.  Both adjoints are \(f_\eta^*Rf_{\eta*}L\boxtimes
g_\eta^*Rg_{\eta*}M\to L\boxtimes M\), the tensor product of the two counits. By the triangle identities of the adjunctions and their compatibility with composition, the maps themselves agree.
\end{proof}

% text: prose-E-12
\phantomsection\label{bp32}%
A complex supported on finitely many closed points is determined by its stalks at these points, and we can read off these stalks from global sections.

% object: lemma-finite-support-global-sections (blueprint/chap-Specialization-and-contracting-correspondences/sec-Kunneth-formula-for-nearby-cycles/lemma-finite-support-global-sections/_section.md)
\phantomsection\label{bp33}%
\begin{bpstatementit}{Lemma~E.11}{ (Complexes with finite support)}
Let \(Y\) be a scheme of finite type over a separably closed field, and let \(Q\in D^b_c(Y_{\acute et},A)\) be supported on a finite set of closed points. If \(R\Gamma(Y,Q)=0\), then \(Q=0\).
\end{bpstatementit}

% proof: proof/lemma-finite-support-global-sections
\begin{proof}[{Proof}]\phantomsection\label{bp34}%
Let \(i_T:T\hookrightarrow Y\) be the finite closed reduced subscheme with underlying set the support of \(Q\), so that \(Q\simeq i_{T*}i_T^*Q\). Since the residue fields of the points of \(T\) are finite and purely inseparable over the separably closed base field, \(T\) is étale locally a finite disjoint union of geometric points, and \(R\Gamma(Y,Q)=\bigoplus_{t\in T}Q_{\bar t}\). Every stalk of \(Q\) therefore vanishes, and \(Q=0\).
\end{proof}

% text: prose-E-13
\phantomsection\label{bp35}%
To restrict to coordinate lines we need coordinates in which the special fibre is quasi-finite over an affine space of the right dimension.

% object: lemma-quasi-finite-coordinates (blueprint/chap-Specialization-and-contracting-correspondences/sec-Kunneth-formula-for-nearby-cycles/lemma-quasi-finite-coordinates/_section.md)
\phantomsection\label{bp36}%
\begin{bpstatementit}{Lemma~E.12}{ (Quasi-finite coordinates)}
Let \(X\) be flat, separated and of finite type over \(S\), and let \(x\) be a closed point of \(X_s\). Put \(q=\dim_xX_s\). Then \(q\leq\dim X_{\bar\eta}\), and there are an open neighbourhood \(V\) of \(x\) in \(X\) and a quasi-finite morphism \(b:V\to\mathbf A^q_S\). In particular \(\dim V_{\bar\eta}\leq q\).
\end{bpstatementit}

% proof: proof/lemma-quasi-finite-coordinates
\begin{proof}[{Proof}]\phantomsection\label{bp37}%
The inequality \(q\leq\dim X_{\bar\eta}\) holds by flatness. Choose a system of parameters \(f_1,\ldots,f_q\) of the local ring \(\mathcal O_{X_s,x}\) of the special fibre, defined on an affine neighbourhood, and lift them to functions on the corresponding affine open of \(X\). The resulting map \(b:X\to\mathbf A^q_S\) is quasi-finite at \(x\), because its fibre through \(x\) is the fibre of \(b_s=(f_1,\ldots,f_q)\), in which \(x\) is isolated. Let \(V\) be the open locus where \(b\) is quasi-finite; as \(b_\eta\) is quasi-finite on \(V_\eta\), the generic fibre of \(V\) has dimension at most \(q\).
\end{proof}

% text: prose-E-14
\phantomsection\label{bp38}%
In the slicing step we compare \(\Pi\) with the Künneth map over the trait \(S'\) of (E.2). Let \(X,Y\) be separated of finite type over \(S\), let \(b:X\to\mathbf A^1_S\) be a morphism, and put

\[
X'=X\times_{\mathbf A^1_S}S',\qquad Y'=Y\times_S S',\qquad
L'=L|_{X'_{\eta'}},\quad M'=M|_{Y'_{\eta'}}.
\]

% object: lemma-kunneth-map-slice (blueprint/chap-Specialization-and-contracting-correspondences/sec-Kunneth-formula-for-nearby-cycles/lemma-kunneth-map-slice/_section.md)
\phantomsection\label{bp39}%
\begin{bpstatementit}{Lemma~E.13}{ (The Künneth map on a slice)}
In the situation just described, let \(L\in D^b_{\mathrm{ctf}}(X_\eta,A)\) and \(M\in D^b_{\mathrm{ctf}}(Y_\eta,A)\). Identify the restrictions over \(s'\) of the nearby cycles of \(L\), \(M\) and \(L\boxtimes M\) by

\[
\begin{aligned}
R\Psi_{X/S}L|_{s'}&\simeq(R\Psi_{X'/S'}L')^P,\\
R\Psi_{Y/S}M|_{s'}&\simeq R\Psi_{Y'/S'}M',\\
R\Psi_{(X\times Y)/S}(L\boxtimes M)|_{s'}
&\simeq
(R\Psi_{(X'\times Y')/S'}(L'\boxtimes M'))^P,
\end{aligned}
\]

given by (E.2) for \(b\), by invariance of nearby cycles under change of trait, and by (E.2) for \(b\circ\operatorname{pr}_1\). Then the restriction of \(\Pi_{L,M}\) over \(s'\) is \((\Pi_{L',M'})^P\). In particular, if \(\Pi_{L',M'}\) is an isomorphism, then so is the restriction of \(\Pi_{L,M}\) over \(s'\).
\end{bpstatementit}

% proof: proof/lemma-kunneth-map-slice
\begin{proof}[{Proof}]\phantomsection\label{bp40}%
In the middle line \(P\) acts trivially, because the action on this complex of nearby cycles is the restriction of an action of \(I\) and \(P\) is the kernel of \(I'\to I\). By (E.3), therefore, \((-)^P\) of the source of \(\Pi_{L',M'}\) is \((R\Psi_{X'/S'}L')^P\boxtimes R\Psi_{Y'/S'}M'\), the restriction of the source of \(\Pi_{L,M}\). By the compatibility of (E.2) with the product maps obtained from \(\bar i^*R\bar j_*\), the map \((\Pi_{L',M'})^P\) is the restriction of \(\Pi_{L,M}\).
\end{proof}

% proof: proof/proposition-external-kunneth-formula-nearby-cycles
\begin{proof}[{Proof of Proposition E.9}]\phantomsection\label{bp41}%
We argue by induction on \(n=\dim(X_{\bar\eta}\times_{\bar\eta}Y_{\bar\eta})\), simultaneously for all such traits.  Before slicing we need two reductions: global sections require proper families, and bounding fibre dimensions after slicing requires flatness, so we arrange both.  For properness, note first that the assertion is local on \(X_s\times_kY_s\).  We may therefore pass to affine neighbourhoods of the two factors, take projective closures over \(S\) and extend the coefficients by zero on their generic fibres. By the base change isomorphism for nearby cycles along open immersions, the restriction of the new map is the original map, so we may assume that \(X\) and \(Y\) are projective.  For flatness, the base change isomorphism for the closed immersions of the closures of their generic fibres, together with invariance under nilpotent thickenings, lets us also assume that they are reduced with dense generic fibres, hence flat over the trait.  Moreover, none of these reductions increases \(n\).

Let \(Q\) be the cone of (E.4). For proper families, global sections of nearby cycles on the special fibre are cohomology of the geometric generic fibre, so by the base change isomorphism for the proper structure maps and Lemma \hyperref[bp30]{E.10}, the map of global complexes induced by \(\Pi_{L,M}\) is the ordinary Künneth isomorphism

\[
R\Gamma(X_{\bar\eta},L)\otimes_A^{\mathbf L}
R\Gamma(Y_{\bar\eta},M)
\xrightarrow{\sim}
R\Gamma(X_{\bar\eta}\times Y_{\bar\eta},L\boxtimes M),
\]

and therefore

\[
R\Gamma(X_s\times_kY_s,Q)=0.
\tag{E.5}
\]

By Lemma \hyperref[bp33]{E.11}, \(Q=0\) as soon as \(Q\) has finite support.  When \(n=0\), the special fibres are zero-dimensional by flatness, so the support is finite. Let \(n>0\). By Lemma \hyperref[bp36]{E.12}, \(X_s\times_kY_s\) is covered by finitely many products \(V_s\times W_s\) of open subschemes with quasi-finite maps \(b:V\to\mathbf A^q_S\) and \(c:W\to\mathbf A^r_S\), where \(q+r\leq n\) and \(\dim V_{\bar\eta}\leq q\), \(\dim W_{\bar\eta}\leq r\). The restriction of \(Q\) to \(V_s\times W_s\) is the cone of \(\Pi_{L|_V,M|_W}\). If \(q+r<n\), it vanishes by the induction hypothesis. By symmetry we may therefore assume \(q+r=n\) and \(q>0\), and compose \(b\) with one of the coordinate projections \(\mathbf A^q_S\to\mathbf A^1_S\). In the notation of Lemma \hyperref[bp39]{E.13} for \(V\), \(W\) and this map, the map \(V'\to\mathbf A^q_S\times_{\mathbf A^1_S}S'=\mathbf
A^{q-1}_{S'}\) is a base change of \(b\), hence quasi-finite, and therefore

\[
\dim V'_{\bar\eta'}\leq q-1,\qquad
\dim W'_{\bar\eta'}\leq r.
\]

By the induction hypothesis \(\Pi_{L',M'}\) is an isomorphism. By Lemma \hyperref[bp39]{E.13}, the restriction of \(Q\) over \(s'\) is therefore zero. In other words, the support of \(Q\) has no point over the generic point of this coordinate line in the special fibre. Its image in that line is constructible and does not contain the generic point, so it is a finite set of closed points.

We repeat this argument for each of the \(q\) coordinates of \(b\). After interchanging the factors, we do the same for each of the \(r\) coordinates of \(c\); the analogue of Lemma \hyperref[bp39]{E.13} with the roles of \(X\) and \(Y\) exchanged holds by symmetry of \(\boxtimes\). The support of \(Q\) on \(V_s\times W_s\) is then contained in the inverse image under \(b_s\times c_s\) of a finite set. This inverse image is finite, because \(b_s\times c_s\) is quasi-finite. With finite support, (E.5) and Lemma \hyperref[bp33]{E.11} show that \(Q=0\).
\end{proof}


% object: sec-Base-change-for-nearby-cycles (blueprint/chap-Specialization-and-contracting-correspondences/sec-Base-change-for-nearby-cycles/_section.md), content hidden
\subsection*{\texorpdfstring{E.3\quad Base change for nearby cycles}{E.3 Base change for nearby cycles}}
\phantomsection\label{bp42}%
\addcontentsline{toc}{subsection}{\texorpdfstring{E.3\quad Base change for nearby cycles}{E.3 Base change for nearby cycles}}
% text: prose-E-15
\phantomsection\label{bp43}%
Let \(c\) be a correspondence over the trait \(S\) and \(u\) a cohomological correspondence on its generic fibre. Nearby cycles turn \(u\) into a map between \(\Psi\) of \(c_2^*K\) and \(\Psi\) of \(c_1^!K\). To obtain a cohomological correspondence \(u_s\) on the special fibre, we must move \(\Psi\) past \(c_2^*\) and \(c_1^!\). This needs base change maps between nearby cycles and the operations \(f^*\), \(Rf_*\), \(Rf_!\) and \(f^!\). We make these maps explicit here. We also prove the identities between them that Proposition \hyperref[bp69]{E.22} uses to show that the trace of \(u_s\) is the specialization of the trace of \(u\).

We write \(\Psi_X=R\Psi_{X/S}\) for a separated \(S\)-scheme \(X\) of finite type. For a morphism \(f:V\to W\) of separated \(S\)-schemes of finite type, put

\[
\begin{array}{rcl}
\alpha_f:f_s^*\Psi_W&\longrightarrow&\Psi_Vf_\eta^*,\\
\chi_f:\Psi_WRf_{\eta*}&\longrightarrow&Rf_{s*}\Psi_V,\\
\nu_f:Rf_{s!}\Psi_V&\longrightarrow&\Psi_WRf_{\eta!},\\
\beta_f:\Psi_Vf_\eta^!&\longrightarrow&f_s^!\Psi_W.
\end{array}
\]

The map \(\alpha_f\) is the base change map for pullback of \(\bar i^*R\bar j_*\), and \(\chi_f\) is its adjoint mate for \(f_s^*\dashv Rf_{s*}\) and \(f_\eta^*\dashv Rf_{\eta*}\).

To define \(\nu_f\), we choose a relative Nagata compactification \(f=ph\), with \(h:V\hookrightarrow\overline V\) open and \(p\) proper. We define \(\nu_h:h_{s!}\Psi_VM\to \Psi_{\overline V}h_{\eta!}M\) using the adjunction for the open immersion, by requiring that its adjoint be

\[
\Psi_VM=\Psi_Vh_\eta^*h_{\eta!}M
\xrightarrow{\alpha_h^{-1}}h_s^*\Psi_{\overline V}h_{\eta!}M.
\]

Since \(p\) is proper, \(\chi_p\) is invertible, and we put \(\nu_f=\chi_p^{-1}\,Rp_{s*}(\nu_h)\), using \(Rf_{s!}=Rp_{s*}h_{s!}\) and \(Rf_{\eta!}=Rp_{\eta*}h_{\eta!}\). The map \(\nu_f\) does not depend on the compactification. Indeed, if \(q\) is a common proper refinement of two compactifications, proper base change identifies \(Rq_*h'_!\) with \(h_!\) on both fibres. After taking adjoints for the open immersions, compatibility with this identification reduces to the composition rule for \(\alpha\) and the triangle identity of the ordinary adjunction. Common refinements also show that \(\nu\) is compatible with composition.

We can also describe \(\nu_f\) on the total families. Write \(i_V\) and \(j_V\) for the inclusions \(\bar i\) and \(\bar j\) of the special and geometric generic fibres of \(V\), so that \(\Psi_V=i_V^*Rj_{V*}\). By proper base change, \(i_W^*Rf_!\simeq Rf_{s!}i_V^*\) and \(j_W^*Rf_!\simeq Rf_{\eta!}j_V^*\). Consider the map \(Rf_!Rj_{V*}\to Rj_{W*}Rf_{\eta!}\) whose \(j_W^*\)-adjoint is

\[
j_W^*Rf_!Rj_{V*}\simeq Rf_{\eta!}j_V^*Rj_{V*}
\xrightarrow{Rf_{\eta!}\varepsilon_{j_V}}Rf_{\eta!}
\]

Its restriction by \(i_W^*\) is a map \(Rf_{s!}\Psi_V\to \Psi_WRf_{\eta!}\), and we claim that this map is \(\nu_f\). As both constructions are compatible with composition, it suffices to check this for \(p\) and for \(h\). For proper \(p\) the restricted map is \(\chi_p^{-1}\), by proper base change for \(i_W^*\). For the open immersion \(h\) we have \(h_\eta^*h_{\eta!}=1\), and the \(h_s^*\)-adjoint of the restricted map is \(\alpha_h^{-1}\), as in the definition of \(\nu_h\).

The map \(\beta_f(M)\) is the adjoint, for \(Rf_{s!}\dashv f_s^!\), of

\[
Rf_{s!}\Psi_Vf_\eta^!M\xrightarrow{\nu_f}
\Psi_WRf_{\eta!}f_\eta^!M
\xrightarrow{\Psi_W\varepsilon_{f_\eta}}\Psi_WM.
\tag{E.6}
\]

For proper \(f\), the maps \(\chi_f\) and \(\nu_f\) are inverse isomorphisms after the identification \(Rf_*=Rf_!\).

The lax monoidal map \(\mu_Y:\Psi_YM\otimes^{\mathbf L} \Psi_YN\to \Psi_Y(M\otimes^{\mathbf L} N)\) of \(R\bar j_*\) is adjoint to the tensor product of the two counits \(j^*Rj_*\to1\). It is associative, symmetric and unital, with unit \(\sigma_Y:A_{Y_s}\to \Psi_YA_{Y_\eta}\), and \(\alpha_f\) is monoidal. For \(f:V\to W\) and \(g:W\to Y\), we have

\[
\alpha_{gf}=\alpha_f(g_\eta^*)f_s^*\alpha_g,
\qquad \sigma_V=\alpha_f(A)f_s^*\sigma_W,
\qquad\mu(\sigma\otimes1)=1.
\tag{E.7}
\]

For the structure morphism \(a_Y:Y\to S\), set \(\theta_Y=\beta_{a_Y}(A):\Psi_Y\omega_{Y_\eta}\to\omega_{Y_s}\), using \(\Psi_SA=A\).

The proof of Proposition \hyperref[bp69]{E.22} moves the base change maps past the projection formula, composition, external duality and the cartesian square defining the fixed point scheme. The following four lemmas give the identities it needs. Each is proved by taking adjoints until both sides become the same composite of base change maps and counits, and then using the composition laws for the base change maps and for the counits. For a morphism \(f\) on either fibre, write \(t_f:f^!M\otimes^{\mathbf L} f^*N\to f^!(M\otimes^{\mathbf L} N)\) for the projection-formula mate of the exceptional counit.

% object: lemma-compatibilities-base-change-maps (blueprint/chap-Specialization-and-contracting-correspondences/sec-Base-change-for-nearby-cycles/lemma-compatibilities-base-change-maps/_section.md)
\phantomsection\label{bp44}%
\begin{bpstatementit}{Lemma~E.14}{ (Base change maps and the projection formula)}
Let \(f:V\to W\) be a morphism of separated \(S\)-schemes of finite type, and let \(M,N\in D^b_c(W_\eta,A)\). Then the two maps \(\Psi_Vf_\eta^!M\otimes^{\mathbf L} f_s^*\Psi_WN\to f_s^!\Psi_W(M\otimes^{\mathbf L} N)\)

\[
\beta_f \Psi_V(t_{f_\eta})\mu_V(1\otimes\alpha_f)
=f_s^!\mu_W\ t_{f_s}(\beta_f\otimes1)
\tag{E.8}
\]

are equal.
\end{bpstatementit}

% proof: proof/lemma-compatibilities-base-change-maps
\begin{proof}[{Proof}]\phantomsection\label{bp45}%
The identity follows, after taking \(Rf_{s!}\)-adjoints, from the compatibility of \(\nu_f\) with the projection formula: for \(L\) on \(V_\eta\) and \(N\) on \(W_\eta\), the composite

\[
Rf_{s!}(\Psi_VL\otimes^{\mathbf L} f_s^*\Psi_WN)\xrightarrow{\mu_V(1\otimes\alpha_f)}
Rf_{s!}\Psi_V(L\otimes^{\mathbf L} f_\eta^*N)\xrightarrow{\nu_f}
\Psi_WRf_{\eta!}(L\otimes^{\mathbf L} f_\eta^*N)\simeq \Psi_W(Rf_{\eta!}L\otimes^{\mathbf L} N)
\]

equals the projection formula on the special fibre followed by \(\mu_W(\nu_f\otimes1)\). To check this, we use the description of \(\nu_f\) on the total families and compare \(j_W^*\)-adjoints. On both sides the adjoint is the counit \(j_V^*Rj_{V*}\to1\) tensored with the identity, followed by the projection formula on the generic fibre.

Taking \(Rf_{s!}\)-adjoints in (E.8) and applying the projection formula, the left side becomes the following composite by this compatibility, and the right side becomes it by naturality of \(\mu_W\):

\[
\begin{split}
Rf_{s!}\Psi_Vf_\eta^!M\otimes^{\mathbf L} \Psi_WN
&\xrightarrow{\nu_f\otimes1}
\Psi_W(Rf_{\eta!}f_\eta^!M)\otimes^{\mathbf L} \Psi_WN\\
&\xrightarrow{\mu_W}\Psi_W(Rf_{\eta!}f_\eta^!M\otimes^{\mathbf L} N)
\xrightarrow{\Psi_W(\varepsilon_{f_\eta}\otimes1)}\Psi_W(M\otimes^{\mathbf L} N).
\end{split}
\]

Since the adjoints agree, so do the maps.
\end{proof}

% object: lemma-composition-exceptional-base-change-maps (blueprint/chap-Specialization-and-contracting-correspondences/sec-Base-change-for-nearby-cycles/lemma-composition-exceptional-base-change-maps/_section.md)
\phantomsection\label{bp46}%
\begin{bpstatementit}{Lemma~E.15}{ (Composition of the exceptional base change maps)}
Let \(f:V\to W\) and \(g:W\to Y\) be morphisms of separated \(S\)-schemes of finite type. Then

\[
\beta_{gf}=f_s^!\beta_g\ \beta_f(g_\eta^!),
\qquad\theta_V=f_s^!\theta_W\ \beta_f(\omega_{W_\eta}).
\tag{E.9}
\]
\end{bpstatementit}

% proof: proof/lemma-composition-exceptional-base-change-maps
\begin{proof}[{Proof}]\phantomsection\label{bp47}%
The first identity follows by taking adjoints for the adjunctions \(Rh_{s!}\dashv h_s^!\) (\(h=f,g,gf\)) and using the compatibility of \(\nu\) with composition together with the composition of counits \(\varepsilon_{gf}=\varepsilon_g\,Rg_!\varepsilon_f\). The second is the case \(g=a_W\), applied to \(A\), of the first, since \(a_V=a_Wf\).
\end{proof}

% text: prose-E-16
\phantomsection\label{bp48}%
Let \(p_1,p_2\) be the projections of \(X\times_SX\). On either fibre, let \(\xi:p_2^*\omega_X\to p_1^!A_X\) be the isomorphism of Lemma \bpout{11.7} for \(K=A_X\), and for a complex \(M\) put \(\kappa_M=t_{p_1}(\xi\otimes1):p_2^*\omega_X\otimes^{\mathbf L} p_1^*M\to p_1^!M\), where \(t_{p_1}:p_1^!A_X\otimes^{\mathbf L} p_1^*M\to p_1^!M\) is the projection-formula mate of the exceptional counit.

The next identity says that, under the base change maps, the map \(\kappa_M\) of external duality on the generic fibre corresponds to \(\kappa_{\Psi_XM}\) on the special fibre, up to \(\theta_X\).

% object: lemma-base-change-maps-external-duality (blueprint/chap-Specialization-and-contracting-correspondences/sec-Base-change-for-nearby-cycles/lemma-base-change-maps-external-duality/_section.md)
\phantomsection\label{bp49}%
\begin{bpstatementit}{Lemma~E.16}{ (Base change maps and external duality)}
Let \(X\) be a separated \(S\)-scheme of finite type with projections \(p_1,p_2\) of \(X\times_SX\), and let \(M\in D^b_c(X_\eta,A)\). Then the two maps \(p_{2s}^*\Psi_X\omega_{X_\eta}\otimes^{\mathbf L} p_{1s}^*\Psi_XM\to p_{1s}^!\Psi_XM\)

\[
\beta_{p_1}\Psi(\kappa_M)\mu
(\alpha_{p_2}\otimes\alpha_{p_1})
=\kappa_{\Psi_XM}(p_{2s}^*\theta_X\otimes1)
\tag{E.10}
\]

are equal.
\end{bpstatementit}

% proof: proof/lemma-base-change-maps-external-duality
\begin{proof}[{Proof}]\phantomsection\label{bp50}%
Since \(\kappa_M=t_{p_1}(\xi\otimes1)\), the main point is the case of the constant complex,

\[
\beta_{p_1}\Psi(\xi)\alpha_{p_2}
=p_{1s}^!\sigma_X\,\xi_s\,p_{2s}^*\theta_X,
\tag{C}
\]

and the general case follows from (C) by Lemma \hyperref[bp44]{E.14}. Since the proof uses proper base change, we prove (C) for a proper family and then extend it by compactification.

Suppose first that \(a=a_X:X\to S\) is proper. On the fibre over \(\eta\), let

\[
Q_\eta(L):a_\eta^*Ra_{\eta*}L\xrightarrow{\sim}Rp_{1\eta*}p_{2\eta}^*L
\]

be the proper base change isomorphism, and define \(Q_s\) in the same way. On each fibre, the \(Rp_{1!}\)-adjoint of \(\xi\) is \(Q(\omega_X)^{-1}\) followed by \(a^*\varepsilon_a\). We compare these adjoints using the identity

\[
\nu_{p_1}Rp_{1s*}(\alpha_{p_2})Q_s(\Psi_XL)
=\Psi_X(Q_\eta(L))\alpha_a(Ra_{\eta*}L)\,a_s^*\nu_a(L).
\]

To prove it, we take adjoints for \(p_{1s}^*\dashv Rp_{1s*}\). By the composition rule for \(\alpha\) and \(ap_1=ap_2\), both adjoints are the composite of \(a_s^*\nu_a\), the base change map \(\alpha_{ap_1}=\alpha_{ap_2}\), and the pullback by \(p_{2\eta}\) of the counit \(a_\eta^*Ra_{\eta*}L\to L\). Substituting this identity into the \(Rp_{1s!}\)-adjoint of the left side of (C), and using naturality of \(\alpha_a\) with respect to \(\varepsilon_{a_\eta}\) and \(\alpha_a(A)=\sigma_X\), we find that this adjoint is the composite

\[
a_s^*Ra_{s!}\Psi_X\omega_{X_\eta}
\xrightarrow{a_s^*\nu_a}
a_s^*Ra_{\eta!}\omega_{X_\eta}
\xrightarrow{a_s^*\varepsilon_{a_\eta}}A_{X_s}
\xrightarrow{\sigma_X}\Psi_XA_{X_\eta}.
\]

By (E.6) for \(a\), this composite equals \(\sigma_X a_s^*(\varepsilon_{a_s}Ra_{s!}\theta_X)\), which is the \(Rp_{1s!}\)-adjoint of the right side. This proves (C) for proper \(X\).

If \(X\) is not proper over \(S\), we choose an open immersion \(j:X\hookrightarrow\overline X\) into a proper \(S\)-scheme and restrict both factors of the product from \(\overline X\) to \(X\). Both sides of (C) are compatible with this restriction. For the maps \(\xi\) this follows from the composition rule for exceptional inverse images. For the base change maps it holds because \(j^!=j^*\), the map \(\alpha_j\) is invertible by smooth base change, and \(\beta_j=\alpha_j^{-1}\). So (C) on \(\overline X\) implies it on \(X\). Applying Lemma \hyperref[bp44]{E.14} to \(p_1\) and using \(\mu(\sigma\otimes1)=1\), we deduce (E.10).
\end{proof}

% text: prose-E-17
\phantomsection\label{bp51}%
For the last identity, let \(c:C\to X\times_SX\) be a correspondence of separated \(S\)-schemes of finite type, with fixed scheme \(F\), closed immersion \(i:F\hookrightarrow C\), map \(d:F\to X\), and closed diagonal \(\Delta:X\hookrightarrow X\times_SX\), so that \(F=C\times_{X\times X}X\). Let \(\zeta:c^!R\Delta_*\xrightarrow\sim Ri_*d^!\) be the base change isomorphism for exceptional inverse image and proper direct image in this cartesian square, on either fibre.

% object: lemma-base-change-maps-fixed-point-scheme (blueprint/chap-Specialization-and-contracting-correspondences/sec-Base-change-for-nearby-cycles/lemma-base-change-maps-fixed-point-scheme/_section.md)
\phantomsection\label{bp52}%
\begin{bpstatementit}{Lemma~E.17}{ (Base change maps and the fixed point scheme)}
In the situation just described,

\[
Ri_{s*}\beta_d\ \chi_i\ \Psi_C(\zeta_\eta)
=\zeta_s\ c_s^!(\chi_\Delta)\ \beta_c(R\Delta_{\eta*}).
\tag{E.11}
\]
\end{bpstatementit}

% proof: proof/lemma-base-change-maps-fixed-point-scheme
\begin{proof}[{Proof}]\phantomsection\label{bp53}%
The inverse of the base change map, \(Ri_*d^!\to c^!R\Delta_*\), is adjoint to the explicit map

\[
Rc_!Ri_*d^!=R\Delta_*Rd_!d^!
\xrightarrow{R\Delta_*\varepsilon_d}R\Delta_*.
\]

This identification uses that \(i\) and \(\Delta\) are proper. Substituting it in both sides, taking \(Rc_{s!}\)-adjoints, and replacing the comparison maps \(\chi_i,\chi_\Delta\) for proper maps by the inverses of \(\nu_i,\nu_\Delta\), we find by the composition rules for the base change maps and for the counits that both sides are the base change map for the composite \(ci=\Delta d\), followed by \(\Psi_X\varepsilon_d\).
\end{proof}


% object: sec-Duality-for-nearby-cycles (blueprint/chap-Specialization-and-contracting-correspondences/sec-Duality-for-nearby-cycles/_section.md), content hidden
\subsection*{\texorpdfstring{E.4\quad Duality for nearby cycles}{E.4 Duality for nearby cycles}}
\phantomsection\label{bp54}%
\addcontentsline{toc}{subsection}{\texorpdfstring{E.4\quad Duality for nearby cycles}{E.4 Duality for nearby cycles}}
% text: prose-E-18
\phantomsection\label{bp55}%
Since the second factor of \(K\boxtimes \mathbb DK\) is a Verdier dual, specializing a trace class also requires nearby cycles to commute with Verdier duality, by an explicit map \(\delta_{X,M}\) defined below [\hyperref[bp99]{Ill94}, Theorem 4.2].

The trait \(S\) and the coefficient ring \(A\) are those of Proposition \hyperref[bp28]{E.9}, and \(\mathbb D_{X_{\bar\eta}}\) and \(\mathbb D_{X_s}\) are the Verdier duals relative to the geometric generic point \(\bar\eta\) and the closed point \(s\) of \(S\). We work on the classical étale site and transport the result to the pro-étale site at the end.

We use the base change maps \(\chi_f\), \(\nu_f\), \(\beta_f\), the lax monoidal map \(\mu_X\) and the map \(\theta_X:\Psi_X\omega_{X_\eta}\to\omega_{X_s}\) of the preceding section.  Recall that \(\beta_f\) is the adjoint, for \(Rf_{s!}\dashv f_s^!\), of (E.6) and that \(\nu_f=\chi_f^{-1}\) for proper \(f\). For \(M\in D^b_{\mathrm{ctf}}(X_\eta,A)\), define the morphism \(\delta_{X,M}:\Psi_X\mathbb D_{X_\eta}M\to\mathbb D_{X_s}(\Psi_XM)\) as the tensor–Hom adjoint of the pairing

\[
\Psi_XM\otimes \Psi_X\mathbb D_{X_\eta}M
\xrightarrow{\mu_X}\Psi_X(M\otimes \mathbb D_{X_\eta}M)
\xrightarrow{\Psi_X\operatorname{ev}}\Psi_X\omega_{X_\eta}
\xrightarrow{\theta_X}\omega_{X_s}.
\tag{E.12}
\]

% object: proposition-nearby-cycles-commute-verdier-duality (blueprint/chap-Specialization-and-contracting-correspondences/sec-Duality-for-nearby-cycles/proposition-nearby-cycles-commute-verdier-duality/_section.md)
\phantomsection\label{bp56}%
\begin{bpstatementit}{Proposition~E.18}{ (Nearby cycles commute with Verdier duality)}
For every separated \(S\)-scheme \(X\) of finite type and \(M\in D^b_{\mathrm{ctf}}(X_\eta,A)\), the map

\[
\delta_{X,M}:\Psi_X\mathbb D_{X_\eta}M
\longrightarrow\mathbb D_{X_s}(\Psi_XM)
\]

defined by (E.12) is an isomorphism.
\end{bpstatementit}

% text: prose-E-19
\phantomsection\label{bp57}%
The argument is the one used for Proposition \hyperref[bp28]{E.9}. As there, on a projective \(X\) global duality shows that the cone of \(\delta\) has zero cohomology on \(X_s\). Restriction to the generic point of a coordinate line (E.2) and induction on dimension show that the cone has finite support. For this we need two compatibilities of \(\delta\): with proper direct image, and with restriction to the generic point of a coordinate line. In the second we keep track of the twist \(A(1)[2]\) on both fibres.

% object: lemma-nearby-duality-proper-direct-image (blueprint/chap-Specialization-and-contracting-correspondences/sec-Duality-for-nearby-cycles/lemma-nearby-duality-proper-direct-image/_section.md)
\phantomsection\label{bp58}%
\begin{bpstatementit}{Lemma~E.19}{ (Duality for nearby cycles and proper direct image)}
Let \(f:X\to Y\) be a proper morphism of separated \(S\)-schemes of finite type, and let \(M\in D^b_{\mathrm{ctf}}(X_\eta,A)\). Under the isomorphism \(\chi_f\) and the Verdier duality isomorphisms \(Rf_*\mathbb D\simeq\mathbb DRf_*\) of Chapter \bpout{10} on both fibres,

\[
Rf_{s*}\delta_{X,M}\quad\text{is identified with}\quad
\delta_{Y,Rf_{\eta*}M}.
\tag{E.13}
\]

For an open immersion \(j:U\to X\), the restriction \(j_s^*\delta_{X,M}\) is \(\delta_{U,j_\eta^*M}\) under the base change isomorphism \(\alpha_j\). If \(X\) is proper over \(S\), then under global Verdier duality on both fibres

\[
R\Gamma(X_s,\delta_{X,M})
\quad\text{is identified with the identity of}\quad
R\operatorname{Hom}_A(R\Gamma(X_{\bar\eta},M),A).
\tag{E.14}
\]
\end{bpstatementit}

% proof: proof/lemma-nearby-duality-proper-direct-image
\begin{proof}[{Proof}]\phantomsection\label{bp59}%
Under Verdier duality for the proper map \(f_\eta\), evaluation on \(Rf_{\eta*}M\) is the composite of the lax monoidal map of \(Rf_{\eta*}\), the map \(Rf_{\eta*}\operatorname{ev}_M\), and the trace \(Rf_{\eta*}\omega_{X_\eta}\to\omega_{Y_\eta}\). Hence, after taking tensor–Hom adjoints, both sides of (E.13) are the pairing

\[
\begin{split}
Rf_{s*}\Psi_XM\otimes Rf_{s*}\Psi_X\mathbb D M
&\xrightarrow{\chi_f^{-1}\otimes\chi_f^{-1}}
\Psi_YRf_{\eta*}M\otimes \Psi_YRf_{\eta*}\mathbb DM\xrightarrow{\mu_Y}
\Psi_Y(Rf_{\eta*}M\otimes Rf_{\eta*}\mathbb DM)\\
&\longrightarrow \Psi_YRf_{\eta*}(M\otimes \mathbb DM)
\longrightarrow \Psi_YRf_{\eta*}\omega_{X_\eta}
\longrightarrow \Psi_Y\omega_{Y_\eta}
\xrightarrow{\theta_Y}\omega_{Y_s}.
\end{split}
\]

To identify the pairing obtained from \(Rf_{s*}\delta_{X,M}\) with this one, we move the lax monoidal maps past the base change maps for proper morphisms, and we use the identity \(\theta_X=f_s^!\theta_Y\circ\beta_f(\omega_{Y_\eta})\) of (E.9). The lax monoidal maps are compatible with the base change maps because, after taking ordinary adjoints, both become the tensor product of the two counits for the open immersions of the generic fibres. This proves (E.13).

For an open immersion \(j\), the map \(\alpha_j\) is an isomorphism, \(j^!=j^*\) and \(\beta_j=\alpha_j^{-1}\). The argument above therefore applies with \(\alpha_j\) in place of \(\chi_f^{-1}\).

For (E.14) we apply (E.13) to the structure morphism \(a\) of \(X\). On the trait, \(\Psi_S\) takes the stalk at the geometric generic point, and the pairing (E.12) is ordinary evaluation. So \(\delta_{S,Ra_{\eta*}M}\) is the identity of \(R\operatorname{Hom}_A(R\Gamma(X_{\bar\eta},M),A)\). Global Verdier duality on both fibres now gives (E.14).
\end{proof}

% text: prose-E-20
\phantomsection\label{bp60}%
For the slice we need that taking invariants under the pro-\(p\) group \(P\) of (E.2) commutes with duality.

% object: lemma-duality-invariants-pro-p (blueprint/chap-Specialization-and-contracting-correspondences/sec-Duality-for-nearby-cycles/lemma-duality-invariants-pro-p/_section.md)
\phantomsection\label{bp61}%
\begin{bpstatementit}{Lemma~E.20}{ (Duality and invariants under $P$)}
Let \(Y\) be a scheme of finite type over a separably closed field, and let \(L\in D^b_{\mathrm{ctf}}(Y_{\acute et},A)\) carry a continuous action of \(P\). There is a natural isomorphism

\[
(\mathbb DL)^P\xrightarrow{\sim}\mathbb D(L^P),
\tag{E.15}
\]

under which the evaluation pairing of \(L^P\) is the restriction of the evaluation pairing of \(L\) to the direct summands \(L^P\) and \((\mathbb DL)^P\).
\end{bpstatementit}

% proof: proof/lemma-duality-invariants-pro-p
\begin{proof}[{Proof}]\phantomsection\label{bp62}%
Let \(G\) be a finite quotient of \(P\). Since \(G\) is a \(p\)-group, the order \(|G|\) is invertible in \(A\), and the invariants are the image of the idempotent \(e=|G|^{-1}\sum_{g\in G}g\).  Verdier duality sends this idempotent to the same averaging idempotent for the contragredient action, since inversion permutes the summands. Hence the image of \(\mathbb D(e)\) is the dual of \(\operatorname{im}e\), and the evaluation pairing restricts as stated.  As invariants are exact, continuity gives the assertion for the pro-\(p\) group \(P\).
\end{proof}

% text: prose-E-21
\phantomsection\label{bp63}%
We fix the notation for the slice. Let \(b:X\to\mathbf A^1_S\) be a coordinate map and let \(S'=(\mathbf A^1_S)_{(s')}\) be the trait at the generic point of the special fibre used in (E.2).  Set \(X'=X\times_{\mathbf A^1_S}S'\), \(M'=M|_{X'_{\bar\eta'}}\), and write \(g_s,g_\eta\) for the two pullbacks to the geometric generic point of the coordinate line.  Restriction to the geometric generic point of the coordinate line is compatible with internal Hom, and Lemma \bpout{10.6} and Lemma \bpout{10.7}, applied on each geometric fibre, give, locally on \(X\), isomorphisms

\[
g_\eta^*\mathbb D_{X_{\bar\eta}}M\simeq \mathbb D_{X'_{\bar\eta'}}M'\otimes A(1)[2],
\qquad
g_s^*\mathbb D_{X_s}L\simeq \mathbb D_{X'_{s'}}(g_s^*L)\otimes A(1)[2]
\tag{E.16}
\]

for \(L\in D^b_{\mathrm{ctf}}(X_s,A)\). The twist comes from the dualizing complexes. For a closed embedding over \(\mathbf A^1\) into \(\mathbf A^r_{\mathbf A^1}\), smooth purity gives \(A(r+1)[2r+2]\) before generic restriction and \(A(r)[2r]\) after it. Moreover, the comparison of the complexes with closed support for the embedding preserves evaluation. Thus on \(\mathbf A^r_{\mathbf A^1_S}\) the identifications of the dualizing complexes before and after slicing differ by \(A(1)[2]\) on both fibres. Finally, the group \(P\) acts trivially on \(A(1)[2]\), because roots of unity of order prime to the residue characteristic lift over a strictly henselian trait. So the twist commutes with taking \(P\)-invariants.

% object: lemma-nearby-duality-slice (blueprint/chap-Specialization-and-contracting-correspondences/sec-Duality-for-nearby-cycles/lemma-nearby-duality-slice/_section.md)
\phantomsection\label{bp64}%
\begin{bpstatementit}{Lemma~E.21}{ (Duality for nearby cycles on a slice)}
In the situation just described, suppose that \(X\) is affine, and let \(i:X\hookrightarrow\mathbf A^r_{\mathbf A^1_S}\) be a closed immersion over \(\mathbf A^1_S\); we use the identifications (E.16) defined by \(i\). Under the identifications (E.2), (E.16) and (E.15), we have

\[
g_s^*\delta_{X,M}
\ =\ (\delta_{X',M'})^P\otimes A(1)[2].
\tag{E.17}
\]
\end{bpstatementit}

% proof: proof/lemma-nearby-duality-slice
\begin{proof}[{Proof}]\phantomsection\label{bp65}%
The map \(\delta\) is the adjoint of the pairing (E.12). By Lemma \hyperref[bp61]{E.20} it therefore suffices to compare the three arrows of this pairing with the corresponding arrows on the slice, restricted to \(P\)-invariants.  Since the slicing formula comes from factoring the inclusion of the geometric generic point and composing the adjunctions for the direct images, the lax monoidal map and the evaluation correspond under it to the maps on the slice, restricted to \(P\)-invariants.

For the last arrow \(\theta_X\), we claim that under (E.2) and (E.16), the map \(g_s^*\theta_X\) corresponds to \((\theta_{X'})^P\otimes A(1)[2]\). We check this on a smooth ambient space, where \(\theta\) is the identity by smooth purity, and descend along the closed immersion \(i:X\hookrightarrow\mathbf A^r_{\mathbf A^1_S}\) of the statement, which we abbreviate to \(i:X\hookrightarrow\mathbf A^r\). First, on \(\mathbf A^r\) and on its slice, smooth purity and the base change isomorphism of nearby cycles for smooth pullback identify \(\theta\) with the identity of the dualizing complex of a smooth scheme, and these identifications differ by \(A(1)[2]\). To descend to \(X\), we use the identity (E.9) for \(i\), which gives \(\theta_X=i_s^!\theta_{\mathbf A^r}\circ\beta_i(\omega_{\mathbf A^r_\eta})\). By (E.6), \(\theta_X\) is then the adjoint, for \(Ri_{s*}\dashv i_s^!\), of \(\nu_i\), followed by \(\Psi\) of the exceptional counit \(Ri_{\eta*}i_\eta^!\to1\) and by \(\theta_{\mathbf A^r}\). Now the base change map \(\nu_i\) for the proper map \(i\) and the factorization used for slicing commute with this counit, because their ordinary adjoints are the same composite of the counit for the open immersion of the generic fibre and the counit for \(i\). Since generic restriction is compatible with \(i^!\) and with evaluation, the claim for \(\mathbf A^r\) gives it for \(X\) before taking invariants.

It remains to pass to \(P\)-invariants. For a finite quotient \(G\) of \(P\), the additive functor \(i_{s'}^!\) commutes with the action of \(G\), and so it preserves the image of the averaging idempotent. Thus \(i_{s'}^!\) commutes with \((-)^P\), and the same descent works after taking \(P\)-invariants. Since the shift \([2]\) is even, no additional Koszul sign appears.
\end{proof}

% proof: proof/proposition-nearby-cycles-commute-verdier-duality
\begin{proof}[{Proof of Proposition E.18}]\phantomsection\label{bp66}%
We argue by induction on \(n=\dim X_{\bar\eta}\), simultaneously for all traits \(S\) as above, and show that the cone of \(\delta_{X,M}\) has zero global sections and finite support.

\emph{Reductions.} We may assume that \(X\) is projective, reduced, and flat over \(S\), by reductions none of which increases \(n\). The assertion is local and, by Lemma \hyperref[bp58]{E.19}, compatible with restriction to opens. So we may pass to a projective compactification and extend \(M\) by zero on its generic fibre. Next we use the compatibility with proper direct image (E.13) for the closure of that generic fibre, together with invariance under nilpotent thickenings. This lets us assume that \(X\) is reduced and flat over \(S\).

Put \(Q=\Cone(\delta_{X,M})\). It is bounded and constructible by the finiteness theorem for nearby cycles and ordinary Verdier duality. By (E.14), \(R\Gamma(X_s,Q)=0\).

\emph{Finite support.} By Lemma \hyperref[bp33]{E.11}, it remains to show that the support of \(Q\) is finite. By Lemma \hyperref[bp36]{E.12}, \(X_s\) is covered by finitely many opens \(V_s\) with quasi-finite coordinate maps \(V\to\mathbf A^q_S\), \(q\leq n\), and \(\dim V_{\bar\eta}\leq q\). As the vanishing of \(Q\) is local, we may replace each \(V\) by finitely many affine opens covering it, with the restricted coordinate maps. Then for each coordinate map \(V\to\mathbf A^1_S\) there is a closed immersion of \(V\) into some \(\mathbf A^r_{\mathbf A^1_S}\) over \(\mathbf A^1_S\), as required in Lemma \hyperref[bp64]{E.21}.

We claim that \(Q\) vanishes over the generic point of each coordinate line in the special fibre. The slice \(V'\) has generic fibre of dimension at most \(q-1<n\). Hence the claim follows from the induction hypothesis and Lemma \hyperref[bp64]{E.21}. The support of \(Q\) therefore lies over finitely many closed points of each coordinate line. Since the coordinate map is quasi-finite, the support is finite. This argument also covers \(n=0\), where \(q=0\) and \(V_s\) is finite.

\emph{Pro-étale formulation.} The comparison for constructible complexes with finite coefficients transports both the pairing and this isomorphism to the small pro-étale site.
\end{proof}


% object: sec-Specialization-of-the-trace-class (blueprint/chap-Specialization-and-contracting-correspondences/sec-Specialization-of-the-trace-class/_section.md), content hidden
\subsection*{\texorpdfstring{E.5\quad Specialization of the trace class}{E.5 Specialization of the trace class}}
\phantomsection\label{bp67}%
\addcontentsline{toc}{subsection}{\texorpdfstring{E.5\quad Specialization of the trace class}{E.5 Specialization of the trace class}}
% text: prose-E-22
\phantomsection\label{bp68}%
In a continuous family of maps of a compact manifold, the Lefschetz number does not change. The algebraic counterpart is the following proposition: the trace class of the specialized correspondence \(u_s\) is the specialization \(\operatorname{sp}_F\) of the trace class of \(u\). Here \(\operatorname{sp}_F\) is a map of sections of dualizing complexes on the fixed scheme. A trace class is built from three pieces: the cohomological correspondence \(u\), the external product \(K\boxtimes\mathbb DK\) in which \(u\) becomes a section, and the evaluation map along the diagonal to the dualizing complex. Nearby cycles commute with each piece. They commute with external products by Proposition \hyperref[bp28]{E.9}, with Verdier duality by Proposition \hyperref[bp56]{E.18}, and with the four operations through the base change maps of Section \hyperref[bp42]{E.3}. We prove the proposition by putting these compatibilities together. Theorem \hyperref[bp91]{E.30} uses it to transfer the computation to the normal cone. For this it also needs Corollary \hyperref[bp78]{E.25}: \(\operatorname{sp}_F\) is the identity when the fixed scheme is, up to nilpotents, a constant family \(Y\times S\).

% object: proposition-specialization-trace-class (blueprint/chap-Specialization-and-contracting-correspondences/sec-Specialization-of-the-trace-class/proposition-specialization-trace-class/_section.md)
\phantomsection\label{bp69}%
\begin{bpstatementit}{Proposition~E.22}{ (Specialization of the trace class)}
Let \(c=(c_1,c_2):C\to X\times_SX\) be a correspondence of separated schemes of finite type over the trait \(S\), and let \(F=\Fix(c)\) be its fixed scheme, with \(i=i_c:F\hookrightarrow C\) and \(d:F\to X\). For \(K\in D^b_{\mathrm{ctf}}(X_\eta,A)\) and a morphism \(u:c_{2\eta}^*K\to c_{1\eta}^!K\), put \(K'=\Psi_XK\) and let \(u_s\) be the composite

\[
u_s:c_{2s}^*K'\xrightarrow{\alpha_{c_2}}\Psi_Cc_{2\eta}^*K
\xrightarrow{\Psi_Cu}\Psi_Cc_{1\eta}^!K
\xrightarrow{\beta_{c_1}}c_{1s}^!K'.
\]

For a section \(z:A_{F_\eta}\to\omega_{F_\eta}\) of the dualizing complex, its specialization is

\[
\operatorname{sp}_F(z)=\theta_F \Psi_F(z)\sigma_F:
A_{F_s}\longrightarrow\omega_{F_s}.
\]

Then

\[
\operatorname{Tr}_{c_s}(u_s)
=\operatorname{sp}_F\bigl(\operatorname{Tr}_{c_\eta}(u)\bigr).
\]
\end{bpstatementit}

% text: prose-E-23
\phantomsection\label{bp70}%
Neither properness of the legs \(c_1,c_2\) nor smoothness of \(F\) is assumed. This matters, because we apply the proposition to the deformation of a correspondence to the normal cone. Its generic fibre is the original correspondence, but its fixed scheme, which we describe in Lemma \hyperref[bp86]{E.28} of Section \hyperref[bp80]{E.6}, may be singular and nonreduced.

A trace class has two ingredients besides \(u\): the external duality isomorphism, which turns \(u\) into a section of \(c^!(K\boxtimes\mathbb DK)\), and the evaluation along the diagonal. In the proof we compare both once and for all on \(X\times_SX\), for the universal kernel \(K\boxtimes\mathbb DK\), and then pull the comparison back along \(c\). After taking adjoints, each comparison reduces to the pairing (E.12) that defines \(\delta_K\).

We fix the maps. In the setting of Proposition \hyperref[bp69]{E.22}, let \(P_X=K\boxtimes\mathbb DK\) be the external product of Lemma \bpout{11.7}.  The comparison map for external products is

\[
\Pi_{M,N}=\mu(\alpha_{p_1}\otimes\alpha_{p_2}):
\Psi_XM\boxtimes \Psi_XN\longrightarrow \Psi_{X\times_SX}(M\boxtimes N),
\]

and it is invertible by Proposition \hyperref[bp28]{E.9}. Using it and the isomorphism \(\delta_K=\delta_{X,K}\) of Proposition \hyperref[bp56]{E.18}, we define the isomorphism

\[
\rho:\Psi_{X\times_SX}P_X\xrightarrow{\Pi_{K,\mathbb DK}^{-1}}
K'\boxtimes \Psi_X\mathbb DK\xrightarrow{1\boxtimes\delta_K}K'\boxtimes\mathbb DK'.
\]

Write \(\lambda_K:P_X\xrightarrow\sim
\cHom_A(p_2^*K,p_1^!K)\) for the external duality isomorphism of Lemma \bpout{11.7}, and \(e_K:P_X\to R\Delta_*\omega_X\) for the evaluation along the diagonal (Definition \bpout{11.8}). Let

\[
\mathfrak h:\Psi_{X\times_SX}\cHom_A(p_2^*K,p_1^!K)\longrightarrow
\cHom_A(p_{2s}^*K',p_{1s}^!K')
\]

be the comparison obtained from the tensor–Hom mate of \(\Psi(\operatorname{ev})\mu\) by precomposing with \(\alpha_{p_2}\) and postcomposing with \(\beta_{p_1}\). For the correspondence \(c\), let \(\lambda_{c,K}:c^!P_X\xrightarrow\sim\cHom_A(c_2^*K,c_1^!K)\) be the external duality isomorphism on \(C\), and define \(\mathfrak h_c\) as \(\mathfrak h\), with \(\alpha_{c_2}\) and \(\beta_{c_1}\) in place of \(\alpha_{p_2}\) and \(\beta_{p_1}\).

% object: lemma-nearby-cycles-universal-kernel (blueprint/chap-Specialization-and-contracting-correspondences/sec-Specialization-of-the-trace-class/lemma-nearby-cycles-universal-kernel/_section.md)
\phantomsection\label{bp71}%
\begin{bpstatementit}{Lemma~E.23}{ (Nearby cycles of the universal kernel)}
In the setting above,

\[
\mathfrak h\,\Psi(\lambda_K)=\lambda_{K'}\rho,
\qquad
e_{K'}\rho=R\Delta_{s*}\theta_X\,\chi_\Delta\,\Psi(e_K).
\tag{E.18}
\]
\end{bpstatementit}

% proof: proof/lemma-nearby-cycles-universal-kernel
\begin{proof}[{Proof}]\phantomsection\label{bp72}%
For the first equality, we precompose with \(\Pi_{K,\mathbb DK}\) and take tensor–Hom adjoints. Both maps then have source

\[
(p_{1s}^*K'\otimes^{\mathbf L} p_{2s}^*\Psi_X\mathbb DK)\otimes^{\mathbf L} p_{2s}^*K'
\]

and target \(p_{1s}^!K'\). The adjoint defining \(\lambda_K\) is

\[
(p_1^*K\otimes^{\mathbf L} p_2^*\mathbb DK)\otimes^{\mathbf L} p_2^*K
\xrightarrow{1\otimes\widetilde{\operatorname{ev}}_K}
p_1^*K\otimes^{\mathbf L} p_2^*\omega_X
\xrightarrow{\sim}p_2^*\omega_X\otimes^{\mathbf L} p_1^*K
\xrightarrow{\kappa_K}p_1^!K,
\]

where \(\widetilde{\operatorname{ev}}_K\) is \(\mathbb DK\otimes^{\mathbf L} K\xrightarrow{\sim}K\otimes^{\mathbf L} \mathbb DK
\xrightarrow{\operatorname{ev}_K}\omega_X\). In both displays the arrows marked \(\sim\) are the symmetry isomorphisms of the tensor product, and all Koszul signs come from these symmetries. Since \(\alpha\) is monoidal, we can move the ordinary comparison maps past \(\mu\). By (E.10) we know what \(\beta_{p_1}\) does to the last arrow \(\kappa_K\). What remains on the factors \(K'\) and \(\Psi_X\mathbb DK\) is the composite \(\theta_X\Psi_X(\operatorname{ev}_K)\mu_X\). This is the pairing (E.12), that is, evaluation against \(\delta_K\). The result is the adjoint of \(\lambda_{K'}(1\boxtimes\delta_K)\). Since \(\Pi\) is invertible, the first equality follows.

For the second equality, we again precompose with \(\Pi_{K,\mathbb DK}\) and take adjoints for \(\Delta_s^*\dashv R\Delta_{s*}\). Since \(\alpha\) is monoidal and compatible with composition, the map \(\alpha_\Delta\Delta_s^*\Pi_{K,\mathbb DK}\) is \(\mu_X:K'\otimes^{\mathbf L} \Psi_X\mathbb DK\to \Psi_X(K\otimes^{\mathbf L} \mathbb DK)\). Hence both adjoints are the pairing (E.12) defining \(\delta_K\). Since \(\Delta\) is proper, the comparison map for direct image is \(\chi_\Delta\).
\end{proof}

% text: prose-E-24
\phantomsection\label{bp73}%
By pulling the first identity back along \(c\), we compute the section of \(c_s^!P_{X_s}\) that corresponds to \(u_s\).

% object: lemma-section-specialized-correspondence (blueprint/chap-Specialization-and-contracting-correspondences/sec-Specialization-of-the-trace-class/lemma-section-specialized-correspondence/_section.md)
\phantomsection\label{bp74}%
\begin{bpstatementit}{Lemma~E.24}{ (The section defined by the specialized correspondence)}
In the setting above,

\[
\mathfrak h_c\Psi_C(\lambda_{c_\eta,K})
=\lambda_{c_s,K'}\,c_s^!\rho\,\beta_c(P_X).
\tag{E.19}
\]

Consequently, if \(k_u:A_{C_\eta}\to c_\eta^!P_X\) is the section corresponding to \(u\), then the section corresponding to \(u_s\) is

\[
k_{u_s}=c_s^!\rho\,\beta_c\,\Psi_C(k_u)\,\sigma_C.
\]
\end{bpstatementit}

% proof: proof/lemma-section-specialized-correspondence
\begin{proof}[{Proof}]\phantomsection\label{bp75}%
After taking tensor–Hom adjoints, we factor the ordinary comparison map through \(c_2=p_2c\) and the exceptional comparison map through \(c_1=p_1c\), and we apply (E.8) for \(f=c\).  What remains is the first equality of (E.18) pulled back by \(c_s^!\). The formula for \(k_{u_s}\) follows from (E.19), the definition of \(u_s\) and \(\mu(\sigma\otimes1)=1\).
\end{proof}

% proof: proof/proposition-specialization-trace-class
\begin{proof}[{Proof of Proposition E.22}]\phantomsection\label{bp76}%
We compare the trace morphisms on the two fibres and then evaluate them on the sections of Lemma \hyperref[bp74]{E.24}. By Definition \bpout{11.8}, the trace morphism is \(t_{c,K}=\zeta\,c^!e_K\). The second equality of (E.18), together with (E.11) and with (E.9) for \(d\), gives

\[
t_{c_s,K'}c_s^!\rho\,\beta_c
=Ri_{s*}\theta_F\,\chi_i\,\Psi_C(t_{c_\eta,K}).
\]

If we apply this to the section \(k_{u_s}\) and take adjoints for \(i_s^*\dashv Ri_{s*}\), the equality of trace classes reduces to the identity for the units

\[
\chi_i \Psi_C(\eta_{i_\eta})\sigma_C
=Ri_{s*}(\sigma_F)\eta_{i_s}.
\]

Its ordinary adjoint is the identity \(\sigma_F=\alpha_i i_s^*\sigma_C\) of (E.7), followed by the ordinary triangle identity.
\end{proof}

% text: prose-E-25
\phantomsection\label{bp77}%
The proposition is useful only if \(\operatorname{sp}_F\) loses no information. In the situation of the contracting argument the fixed scheme does not move in the family, and then specialization is the identity.

% object: corollary-specialization-when-fixed-point-scheme (blueprint/chap-Specialization-and-contracting-correspondences/sec-Specialization-of-the-trace-class/corollary-specialization-when-fixed-point-scheme/_section.md)
\phantomsection\label{bp78}%
\begin{bpstatementit}{Corollary~E.25}{ (Specialization when the fixed point scheme is a constant family)}
If \(F_{\mathrm{red}}=Y\times S\), with \(Y\) separated and of finite type over \(k\), then the specialization \(\operatorname{sp}_F\) of Proposition \hyperref[bp69]{E.22} is the identity under the identifications for constant families and invariance under nilpotent thickenings. In particular, it is injective.
\end{bpstatementit}

% proof: proof/corollary-specialization-when-fixed-point-scheme
\begin{proof}[{Proof}]\phantomsection\label{bp79}%
By naturality it suffices to show that \(\theta_Y\sigma_Y(\omega_Y)\) is the identity of \(\omega_Y\). For proper \(Y\) this follows from the adjunction defining \(\theta_Y\). We deduce the case of general \(Y\) by a compactification, and the case of nonreduced \(F\) by invariance under nilpotent thickenings.

\emph{Proper \(Y\) and \(F=Y\times S\).} Here \(\theta_Y\sigma_Y(\omega_Y)=\operatorname{id}_{\omega_Y}\). The unit \(\sigma_Y(\omega_Y):\omega_Y\to \Psi_Y\omega_{Y_\eta}\) is an isomorphism by Lemma \hyperref[bp7]{E.2}. For the proper structure map \(a\) of \(Y\), the base change isomorphism for proper maps and its compatibility with \(\sigma\) give

\[
\nu_a\,Ra_!\sigma_Y(\omega_Y)
=\sigma_{\operatorname{Spec}k}(Ra_!\omega_Y)
=\operatorname{id}_{Ra_!\omega_Y}.
\]

By the defining adjunction (E.6), \(\theta_Y\sigma_Y(\omega_Y)\) is therefore the adjoint of the counit \(Ra_!\omega_Y\to A\), which is \(\operatorname{id}_{\omega_Y}\). Naturality of \(\sigma\) then gives \(\theta_Y\Psi_Y(z)\sigma_Y(A)=z\) for every section \(z\) of the dualizing complex.

\emph{General \(Y\).} Choose a proper compactification \(j:Y\hookrightarrow\overline Y\), so that \(\omega_Y=j^*\omega_{\overline Y}\). We compare the two sides of the identity under restriction along \(j\). The units for the constant families correspond by base change along the open immersion, and the map \(\theta_Y\) corresponds to \(j^*\theta_{\overline Y}\) by the composition rule (E.9) together with \(\beta_j=\alpha_j^{-1}\) for the open immersion. Restricting the identity \(\theta_{\overline Y}\sigma_{\overline Y}(\omega_{\overline Y})=1\) therefore gives \(\theta_Y\sigma_Y(\omega_Y)=1\), and by naturality the assertion follows again for all sections of the dualizing complex.

\emph{Nonreduced \(F\).} If only \(F_{\mathrm{red}}=Y\times S\), we replace \(F\) by its reduction. This does not change \(\operatorname{sp}_F\), because invariance under nilpotent thickenings of the total family and of its fibres shows that the six operations, their units and counits, and nearby cycles agree with those of the reduction.
\end{proof}


% object: sec-Contracting-correspondences (blueprint/chap-Specialization-and-contracting-correspondences/sec-Contracting-correspondences/_section.md), content hidden
\subsection*{\texorpdfstring{E.6\quad Contracting correspondences}{E.6 Contracting correspondences}}
\phantomsection\label{bp80}%
\addcontentsline{toc}{subsection}{\texorpdfstring{E.6\quad Contracting correspondences}{E.6 Contracting correspondences}}
% text: prose-E-26
\phantomsection\label{bp81}%
Let \(c=(c_1,c_2):C\to X\times X\) be a correspondence, and let \(Z\hookrightarrow X\) be a closed subscheme. Suppose that the leg \(c_2\), along which we take ordinary pullback, contracts towards \(Z\) faster than \(c_1\), as measured by powers of ideals. Then the trace class near a connected component of the fixed point scheme meeting \(Z\) can be computed by restricting the correspondence to \(Z\). With our convention \(u:c_2^*K\to c_1^!K\) for the legs, the leg that contracts is \(c_2\); in [\hyperref[bp98]{Var05}] it is the first leg.

The contraction becomes geometric on the deformation to the normal cone (Lemma \hyperref[bp82]{E.26} below). Let \(I\) be the ideal of \(Z\) and \(I_i=c_i^{-1}(I)\mathcal O_C\). Under the hypothesis \(I_2^n\subset I_1^{n+1}\), the special fibre of \(c_2\) maps into the zero section, while the inverse image of the zero section under \(c_1\) is the zero section. On the special fibre the trace class therefore only sees \(Z\).

% object: lemma-normal-cones-contracting-ideals (blueprint/chap-Specialization-and-contracting-correspondences/sec-Contracting-correspondences/lemma-normal-cones-contracting-ideals/_section.md)
\phantomsection\label{bp82}%
\begin{bpstatementit}{Lemma~E.26}{ (Normal cones and contracting ideals)}
Let \(c=(c_1,c_2):C\to X\times X\) be a correspondence, let \(Z\hookrightarrow X\) have ideal \(I\), and put

\[
I_1=c_1^{-1}(I)\mathcal O_C,\qquad
I_2=c_2^{-1}(I)\mathcal O_C,\qquad Z'=V(I_1).
\]

Assume that \(I_2\subset I_1\) and \(I_2^n\subset I_1^{n+1}\) for some \(n\geq1\). Over the trait \(T=\operatorname{Spec}k[t]_{(t)}\), whose strict henselization is the trait \(S\) used for nearby cycles, the two legs lift to the deformations to the normal cones,

\[
\widetilde c_i:\mathcal D_{Z'}C\longrightarrow\mathcal D_ZX,
\qquad
\mathcal D_ZX=
\operatorname{Spec}_{X_T}\mathcal O_{X_T}[I/t],
\]

where \(\mathcal D_ZX\) is written over \(T\) and \(\mathcal D_{Z'}C\) is defined in the same way from \(I_1\). Their special fibres \(\gamma_i:N_{Z'}C\to N_ZX\) satisfy:

\begin{enumerate}
\item \(\gamma_1^{-1}(Z)=Z'\) scheme-theoretically, where \(Z\) and \(Z'\) are the zero sections of the two normal cones;
\item the image of \(\gamma_2\) is set-theoretically contained in the zero section \(Z\subset N_ZX\).
\end{enumerate}
\end{bpstatementit}

% proof: proof/lemma-normal-cones-contracting-ideals
\begin{proof}[{Proof}]\phantomsection\label{bp83}%
We read everything off from the graded algebras. The first leg sends \(I\) onto generators of \(I_1\), so it identifies the ideals of the zero sections. The second leg, on the other hand, sends \(I\) into \(I_2\), and by the hypothesis \(I_2^n\subset I_1^{n+1}\) the images of degree-one classes are nilpotent in \(\operatorname{gr}_{I_1}\mathcal O_C\).

The algebra of the normal cone is \(\operatorname{gr}_I\mathcal O_X=\bigoplus_{r\geq0}I^r/I^{r+1}\), and the algebra of the deformation is the subalgebra of \(\mathcal O_{X_T}[t^{-1}]\) generated by \(\mathcal O_{X_T}\) and the fractions \(a/t\), \(a\in I\). Since \(c_i^{-1}(I)\mathcal O_C\subset I_1\) for both legs (for \(i=2\) this is the hypothesis \(I_2\subset I_1\)), the maps on generic fibres preserve these subalgebras. This gives the lifts.

The zero section of \(N_ZX\) is defined by the ideal of elements of positive degree of \(\operatorname{gr}_I\mathcal O_X\). Since \(I_1\) is the ideal generated by \(c_1^{-1}(I)\), the images of elements of degree one under the first leg generate \(I_1/I_1^2\) as an \(\mathcal O_{Z'}\)-module, and hence they generate the ideal of elements of positive degree of \(\operatorname{gr}_{I_1}\mathcal O_C\). Therefore the quotient by the pullback of the ideal of the zero section is \(\mathcal O_C/I_1=\mathcal O_{Z'}\), which proves (1).

Under the second leg, a product of \(n\) images of elements of degree one is the class of an element of \(I_2^n\) in \(I_1^n/I_1^{n+1}\), and this class is zero since \(I_2^n\subset I_1^{n+1}\). The pullback of the ideal of the zero section is therefore nilpotent. Its vanishing set is all of \(N_{Z'}C\), which means that \((N_{Z'}C)_{\mathrm{red}}\to N_ZX\) factors through the zero section. This proves (2).
\end{proof}

% object: example-contracting-squaring (blueprint/chap-Specialization-and-contracting-correspondences/sec-Contracting-correspondences/example-contracting-squaring/_section.md)
\phantomsection\label{bp84}%
\begin{bpstatement}{Example~E.27}{ (Squaring on the line)}
Let \(X=C=\mathbf A^1=\operatorname{Spec}k[x]\), let \(Z\) be the origin, with \(I=(x)\), and let \(c_1=\mathrm{id}\) and \(c_2(x)=x^2\), so that \(c_2^*x=x^2\). Then \(I_1=(x)\) and \(I_2=(x^2)=I_1^2\), so \(c\) contracts towards \(Z\) with \(n=1\). The fixed scheme is \(V(x^2-x)\). It consists of the two reduced points \(0\) and \(1\), and only the first meets \(Z\).

The deformation to the normal cone is \(\operatorname{Spec}k[t]_{(t)}[x,x/t]\). This is the constant family \(\mathbf A^1_T\) in the coordinate \(x'=x/t\). In this coordinate the first leg is the identity and the second is \(x'\mapsto x^2/t=tx'^2\). Hence, on the special fibre \(t=0\), the second leg is the constant map to \(0\), as (2) of the lemma predicts, and the deformed fixed scheme there is the point \(x'=0\). On the generic fibre, by contrast, there are two fixed points: one is \(x'=0\), the constant family of the origin, and the other is \(x'=1/t\), the point \(x=1\). This second point leaves every bounded region of the \(x'\)-line as \(t\to0\), and so it does not reach the special fibre.

Corollary \hyperref[bp95]{E.31} below therefore gives \(\operatorname{LT}_0(u)=\operatorname{tr}_A(u_{\bar 0}\mid K_{\bar
0})\) for every \(K\) and \(u\). For instance, the local term is zero if \(K=j_!A\) with \(j:\mathbf G_m\hookrightarrow\mathbf A^1\), as \(K_{\bar 0}=0\). The condition fails for the identity correspondence. There \(I_2=I_1\), the inclusion \(I_1^n\subset I_1^{n+1}\) is false for every \(n\), and the fixed scheme is all of \(X\).
\end{bpstatement}

% text: prose-E-27
\phantomsection\label{bp85}%
As in Proposition \hyperref[bp69]{E.22}, write \(F=\Fix(c)\) for the fixed scheme, and let \(d=c_1|_F=c_2|_F:F\to X\) be the common map. We will show that a connected component of \(F\) that meets \(Z\) lies in \(Z\) set-theoretically, and that the deformed fixed scheme is constant up to nilpotents. Both follow from the hypothesis \(I_2^n\subset I_1^{n+1}\), because the two legs agree on \(F\).

The deformation of \(F\) to the normal cone of \(d^{-1}(Z)\) embeds in the fixed scheme of the deformed correspondence. Put \(I_F=d^{-1}(I)\mathcal O_F\), so that \(d^{-1}(Z)=V(I_F)\). Since \(d=c_1|_F\), the ideal \(I_F^r\) is the image of \(I_1^r\) in \(\mathcal O_F\) for every \(r\). Hence the closed immersion \(F\hookrightarrow C\) induces a surjection of deformation algebras

\[
\mathcal O_{C_T}[I_1/t]\twoheadrightarrow
\mathcal O_{F_T}[I_F/t],
\]

that is, a closed immersion \(\mathcal D_{d^{-1}(Z)}F\hookrightarrow\mathcal D_{Z'}C\). Its two composites with \(\widetilde c_1\) and \(\widetilde c_2\) are the maps induced by \(c_1|_F\) and \(c_2|_F\), which are equal. So it factors through the fixed scheme \(\widetilde F=\Fix(\widetilde c)\) of \(\widetilde c=(\widetilde c_1,\widetilde c_2)\). We write

\[
\iota_F:\mathcal D_{d^{-1}(Z)}F\longrightarrow\widetilde F
\tag{E.20}
\]

for the resulting closed immersion.

% object: lemma-fixed-point-scheme-deformation (blueprint/chap-Specialization-and-contracting-correspondences/sec-Contracting-correspondences/lemma-fixed-point-scheme-deformation/_section.md)
\phantomsection\label{bp86}%
\begin{bpstatementit}{Lemma~E.28}{ (The fixed point scheme of the deformation)}
In the situation of Lemma \hyperref[bp82]{E.26}, suppose that \(F\) is connected and that \(d(F)\) meets \(Z\). Then \(d^{-1}(Z)_{\mathrm{red}}=F_{\mathrm{red}}\), the reduction of \(\mathcal D_{d^{-1}(Z)}F\) is \((F_T)_{\mathrm{red}}\), and the map \((\iota_F)_{\mathrm{red}}:(F_T)_{\mathrm{red}}\to\widetilde F_{\mathrm{red}}\) induced by (E.20) is an isomorphism.
\end{bpstatementit}

% proof: proof/lemma-fixed-point-scheme-deformation
\begin{proof}[{Proof}]\phantomsection\label{bp87}%
Restricting \(I_2^n\subset I_1^{n+1}\) to \(F\), where the two legs agree, gives \(I_F^n\subset I_F^{n+1}\). The reverse inclusion always holds, so the coherent ideal \(M=I_F^n\) satisfies \(M=I_FM\). Hence, at a point of \(V(I_F)\), Nakayama's lemma shows that \(M=0\) in the local ring, and hence on a neighbourhood of the point, whereas at a point outside \(V(I_F)\) we have \(M=\mathcal O_F\). The closed set \(V(I_F)\) is therefore also open. It is nonempty by assumption, and \(F\) is connected, so \(V(I_F)=F\) and \(I_F^n=0\). This proves the first assertion. As \(I_F\) is nilpotent, each generator \(h/t\) of \(\mathcal O_{F_T}[I_F/t]\) over \(\mathcal O_{F_T}\) is nilpotent. Hence the reduction of \(\mathcal D_{d^{-1}(Z)}F\) is \((F_T)_{\mathrm{red}}\).

It remains to see that \(\widetilde F\) has no further points. On the generic fibre the deformed correspondence is \(c\) itself, so the generic fibre of \(\widetilde F\) is that of \(F_T\). Let \(x\) be a point of the special fibre of \(\widetilde F\). Then \(\gamma_2(x)\) lies in the zero section by Lemma lemma-normal-cones-contracting-ideals, and so does \(\gamma_1(x)=\gamma_2(x)\). By Lemma lemma-normal-cones-contracting-ideals, the point \(x\) lies on the zero section \(Z'\subset N_{Z'}C\). On \(Z'\) the two maps are the restrictions of \(c_1\) and \(c_2\). Their equalizer is \(F\cap Z'\), whose underlying space is \(F\) by the first assertion. Thus every fixed point in the special fibre lies in the closed subscheme above, and the closed immersion of reduced schemes is surjective. Its ideal has zero radical and is therefore zero, so the immersion is an isomorphism.
\end{proof}

% text: prose-E-28
\phantomsection\label{bp88}%
Up to nilpotents, then, the fixed scheme of the deformed correspondence is the constant family \(F_T\). This is why the specialization of trace classes can be inverted in the proof of Theorem \hyperref[bp91]{E.30}. The hypothesis that \(F\) is connected and meets \(Z\) is needed. For instance, in Example \hyperref[bp84]{E.27}, the component \(F=\{0\}\) gives the constant family \(x'=0\), but the fixed point \(x=1\) does not meet \(Z\), and there the deformation gives only the generic-fibre point \(x'=1/t\), not the constant family \(\{1\}_T\).

% object: remark-reduction-connected-fixed-point-scheme (blueprint/chap-Specialization-and-contracting-correspondences/sec-Contracting-correspondences/remark-reduction-connected-fixed-point-scheme/_section.md)
\phantomsection\label{bp89}%
\begin{bpstatement}{Remark~E.29}{ (Reduction to a connected fixed point scheme)}
The hypothesis in the preceding lemma that \(F\) is connected is no loss of generality. Indeed, let \(F_0\subset \Fix(c)\) be a connected component that is open and closed. If we replace \(C\) by an open neighbourhood of \(F_0\) on which the ideal inclusions hold, and then remove \(\Fix(c)\setminus F_0\), we obtain an open subscheme of \(C\), with the same \(X\) and \(K\), whose fixed scheme is \(F_0\). Restriction of trace classes along this open immersion identifies the trace class of \(u\) on \(F_0\) with the full trace class of the restricted correspondence.
\end{bpstatement}

% text: prose-E-29
\phantomsection\label{bp90}%
The proof of the main theorem combines the two lemmas on the deformation with Verdier's theorem and with the specialization of trace classes. The statement of the main theorem involves the restriction of a correspondence to \(Z\), which we construct first.

Let \(c=(c_1,c_2):C\to X\times X\) be a correspondence, let \(Z\hookrightarrow X\) have ideal \(I\), and suppose that \(c_2^{-1}(I)\mathcal O_C\subset c_1^{-1}(I)\mathcal O_C\). Then \(c_1^{-1}(Z)\subset c_2^{-1}(Z)\). As before Lemma \bpout{11.17}, we put

\[
C_Z=(c_1^{-1}Z)_{\mathrm{red}},\qquad
c_Z=(c_{Z,1},c_{Z,2}):C_Z\longrightarrow Z_{\mathrm{red}}\times Z_{\mathrm{red}},
\]

with \(i:Z_{\mathrm{red}}\hookrightarrow X\) and \(i':C_Z\hookrightarrow C\). For \(K\in D^b_{\mathrm{ctf}}(X,A)\) put \(K_Z=i^*K\). The square formed by \(i'\), \(i\), \(c_1\) and \(c_{Z,1}\) is cartesian up to reduction. Hence invariance under nilpotent thickenings and base change for direct image with compact support give \(i^*Rc_{1!}\simeq Rc_{Z,1!}i'^*\). Since the second leg maps \(C_Z\) into \(Z_{\mathrm{red}}\), we also have \(i'^*c_2^*K\simeq c_{Z,2}^*K_Z\). For a cohomological correspondence \(u:c_2^*K\to c_1^!K\) with adjoint \(\widehat u:Rc_{1!}c_2^*K\to K\) under \(Rc_{1!}\dashv c_1^!\), the \emph{restriction} \(u_Z:c_{Z,2}^*K_Z\to c_{Z,1}^!K_Z\) of \(u\) to \(Z\) is the adjoint under \(Rc_{Z,1!}\dashv c_{Z,1}^!\) of

\[
Rc_{Z,1!}c_{Z,2}^*K_Z
\simeq i^*Rc_{1!}c_2^*K
\xrightarrow{i^*\widehat u}K_Z.
\tag{E.21}
\]

Equivalently, \(u_Z\) is \(i'^*u\) followed by the base change map \(i'^*c_1^!K\to c_{Z,1}^!K_Z\), which in general is not an isomorphism. This is the restriction to \(Z\) that occurs in Lemma \bpout{11.18}. If \(F_0\) is an open and closed subscheme of \(\Fix(c)\), put \(F_{0,Z}=\Fix(c_Z)\times_{\Fix(c)}F_0\), and let \(h_{F_0}:F_{0,Z}\to F_0\) be the projection.

% object: theorem-contracting-correspondences-local-terms (blueprint/chap-Specialization-and-contracting-correspondences/sec-Contracting-correspondences/theorem-contracting-correspondences-local-terms/_section.md)
\phantomsection\label{bp91}%
\begin{bpstatementit}{Theorem~E.30}{ (Contracting correspondences and local terms)}
Let \(c=(c_1,c_2):C\to X\times X\) be a correspondence, and let \(Z\hookrightarrow X\) have ideal \(I\). Let \(F_0\) be a connected component of \(\Fix(c)\) that is open and closed and whose image meets \(Z\). Choose an open neighbourhood \(W\subset C\) of \(F_0\) on which

\[
I_2=c_2^{-1}(I)\mathcal O_W\subset
I_1=c_1^{-1}(I)\mathcal O_W,
\qquad I_2^n\subset I_1^{n+1}\quad\text{for some }n\geq1.
\]

Let \(K\in D^b_{\mathrm{ctf}}(X,A)\), let \(u:c_2^*K\to c_1^!K\) be a cohomological correspondence, and let \(u_Z\) be the restriction (E.21) to \(Z\) of \(u|_W\), for the correspondence \(c|_W\), with \(c_Z:C_Z=(W\cap c_1^{-1}Z)_{\mathrm{red}}\to
Z_{\mathrm{red}}\times Z_{\mathrm{red}}\). Then \(h_{F_0}\) is an isomorphism after reduction, and

\[
\operatorname{Tr}_c(u)|_{F_0}
=(h_{F_0})_*\bigl(\operatorname{Tr}_{c_Z}(u_Z)|_{F_{0,Z}}\bigr),
\]

where \((h_{F_0})_*:H^0(F_{0,Z},\omega_{F_{0,Z}})\to H^0(F_0,\omega_{F_0})\) is the induced isomorphism. In particular, if \(F_0\) is proper, then \(\operatorname{LT}_{F_0}(u)=\operatorname{LT}_{F_{0,Z}}(u_Z)\).
\end{bpstatementit}

% text: prose-E-30
\phantomsection\label{bp92}%
For \(Z\) a point, the theorem specializes to Corollary \hyperref[bp95]{E.31}. Under its hypotheses the local term at \(F_0\) is the categorical trace of \(u_Z\). For \(c=(\mathrm{id}_X,h)\) this is \(\operatorname{tr}(u_{\bar y}\mid K_{\bar y})\).

% proof: proof/theorem-contracting-correspondences-local-terms
\begin{proof}[{Proof of Theorem E.30}]\phantomsection\label{bp93}%
We split \(K\) by the triangle \(j_!j^*K\to K\to Ri_*i^*K\) and use additivity of trace classes. The piece supported on \(Z\) gives the trace of \(u_Z\). The piece \(j_!j^*K\), on the other hand, has zero trace, which we show by deforming to the normal cone.

\emph{Reductions.} We may assume that \(W=C\), that \(\Fix(c)=F_0\) and that the ideal inclusions hold on \(C\). Indeed, by Remark \hyperref[bp89]{E.29}, restricting \(C\) to \(W\) and removing the other connected components of \(\Fix(c)\) changes neither the trace class on \(F_0\) nor \(u_Z\) near \(F_{0,Z}\). By Lemma \hyperref[bp86]{E.28}, the common map \(\Fix(c)\to X\) maps set-theoretically into \(Z\). The same lemma shows that after base change from \(T\) to \(S\) the reduction of the fixed scheme of the deformed correspondence is \((\Fix(c)\times S)_{\mathrm{red}}\), as Corollary \hyperref[bp78]{E.25} requires.

\emph{The case \(i^*K=0\).} We show that \(\operatorname{Tr}_c(u)=0\). First, the specialized complex vanishes on the zero section,

\[
\operatorname{Sp}_{X,Z}(K)|_Z=0,
\]

because Verdier's theorem on the zero section (Theorem \hyperref[bp24]{E.8}) identifies the restriction of \(\operatorname{Sp}_{X,Z}(K)\) to the zero section with the restriction of \(K\) to \(Z\), and invariance under the nilpotent thickening \(Z_{\mathrm{red}}\to Z\) passes from \(Z\) to \(Z_{\mathrm{red}}\).

Next, consider the leg \(\gamma_2\) of the correspondence on the special fibre, along which we take ordinary pullback. It maps set-theoretically into the zero section by Lemma lemma-normal-cones-contracting-ideals. Invariance under nilpotent thickenings therefore gives \(\gamma_2^*\operatorname{Sp}_{X,Z}(K)=0\). So the specialized morphism has zero source, and its trace is zero.

Finally, we relate this trace to \(\operatorname{Tr}_c(u)\). The generic fibre of the deformed correspondence is \(c\) over \(\eta\). So Proposition \hyperref[bp69]{E.22} identifies the zero trace with the specialization of the pullback of \(\operatorname{Tr}_c(u)\) to the generic fibre. By Corollary \hyperref[bp78]{E.25}, specialization is the identity under the identifications for constant families, and these carry that pullback to \(\operatorname{Tr}_c(u)\). Hence \(\operatorname{Tr}_c(u)=0\).

\emph{General \(K\).} We apply Proposition proposition-additivity-trace-classes to the triangle

\[
j_!j^*K\longrightarrow K\longrightarrow Ri_*i^*K,
\qquad j:X\setminus Z_{\mathrm{red}}\hookrightarrow X,
\]

whose hypothesis \(c_1^{-1}(Z_{\mathrm{red}})\subset c_2^{-1}(Z_{\mathrm{red}})\) follows from \(I_2\subset I_1\). The open piece has zero trace by the case just treated, while the closed graded piece is the proper pushforward of \(u_Z\) along \(i\) and \(i'\), by Lemma \bpout{11.18}, and Theorem \bpout{11.13} identifies its trace with the pushforward of the trace class of \(u_Z\). To compare with \(F_0\), note that the map \(h_{F_0}\) is an isomorphism after reduction, because \(\Fix(c)\) lies in \(c_1^{-1}(Z)\) set-theoretically by Lemma \hyperref[bp86]{E.28}. This proves the displayed formula. If \(F_0\) is proper, then so is \(F_{0,Z}\), and pushing both sides forward to \(\operatorname{Spec}k\) and composing with the trace map of Definition \bpout{11.9} gives the equality of local terms.
\end{proof}

% text: prose-E-31
\phantomsection\label{bp94}%
If \(Z\) is a point, the restricted correspondence is a correspondence of the point, and the theorem computes the local term as a trace on a stalk.

% object: corollary-contracting-local-terms-point (blueprint/chap-Specialization-and-contracting-correspondences/sec-Contracting-correspondences/corollary-contracting-local-terms-point/_section.md)
\phantomsection\label{bp95}%
\begin{bpstatementit}{Corollary~E.31}{ (Contracting local terms at a point)}
Let \(y\in X(k)\), and suppose that \(c_1^{-1}(y)\) is a single reduced point \(z\), with \(c_2(z)=y\). Suppose that \(F_0=\{z\}\) is an isolated reduced component of \(\Fix(c)\), open and closed, and that on a neighbourhood of \(F_0\) the ideal \(I_y\) satisfies

\[
c_2^{-1}(I_y)\mathcal O_C\subset c_1^{-1}(I_y)\mathcal O_C,
\qquad
\bigl(c_2^{-1}(I_y)\mathcal O_C\bigr)^n
\subset\bigl(c_1^{-1}(I_y)\mathcal O_C\bigr)^{n+1}
\quad\text{for some }n\geq1.
\]

Then for every \(K\in D^b_{\mathrm{ctf}}(X,A)\) and every morphism \(u:c_2^*K\to c_1^!K\), the local term at \(F_0\) is the categorical trace of the restriction \(u_Z\) to \(Z=\{y\}\) of \(u\) restricted to that neighbourhood, defined in (E.21). In particular, for \(c=(\operatorname{id}_X,h)\),

\[
\operatorname{LT}_y(u)
=\operatorname{tr}_A(u_{\bar y}\mid K_{\bar y}).
\]
\end{bpstatementit}

% text: prose-E-32
\phantomsection\label{bp96}%
Here \(X\) may be singular.

% proof: proof/corollary-contracting-local-terms-point
\begin{proof}[{Proof of Corollary E.31}]\phantomsection\label{bp97}%
We apply Theorem \hyperref[bp91]{E.30} with \(Z=\{y\}\) and \(W\) the neighbourhood of the statement; \(F_0\) meets \(Z\) because \(d(z)=c_1(z)=y\). Here \(C_Z\) is the reduced point \(z\) and \(F_{0,Z}=F_0\), so that \(c_Z\) is the identity correspondence of the point, with coefficient \(K_{\bar y}\), and by Definition \bpout{11.9} its local term is the categorical trace. If the first leg is the identity, the base change map in the construction of \(u_Z\) is the identity, and hence the restricted morphism is the stalk map \(u_{\bar y}\).
\end{proof}



\clearpage\phantomsection\hypertarget{bp-outside}{}%
\section*{Outside this export}
\addcontentsline{toc}{section}{Outside this export}
References marked \ensuremath{{}^\ast} are to parts of the blueprint this document leaves out:
\begin{description}
\item[Theorem 2.43] Comparison for constructible sheaves, in Ⅰ The six operations (\bpslug{theorem-{}\allowbreak{}comparison-{}\allowbreak{}constructible-{}\allowbreak{}sheaves})
\item[Corollary 2.44] Finite Tor amplitude under comparison, in Ⅰ The six operations (\bpslug{corollary-{}\allowbreak{}finite-{}\allowbreak{}tor-{}\allowbreak{}amplitude-{}\allowbreak{}under-{}\allowbreak{}comparison})
\item[Proposition 8.18] Ordinary external Künneth, in Ⅰ The six operations (\bpslug{proposition-{}\allowbreak{}ordinary-{}\allowbreak{}external-{}\allowbreak{}kunneth})
\item[Chapter 10] Verdier and Poincaré duality, in Ⅱ Duality and the trace formula (\bpslug{chap-{}\allowbreak{}Verdier-{}\allowbreak{}and-{}\allowbreak{}Poincare-{}\allowbreak{}duality})
\item[Lemma 10.6] Restriction to a geometric generic fibre, in Ⅱ Duality and the trace formula (\bpslug{lemma-{}\allowbreak{}restriction-{}\allowbreak{}geometric-{}\allowbreak{}generic-{}\allowbreak{}fibre})
\item[Lemma 10.7] Double duals on the geometric generic fibre, in Ⅱ Duality and the trace formula (\bpslug{lemma-{}\allowbreak{}double-{}\allowbreak{}dual-{}\allowbreak{}geometric-{}\allowbreak{}generic-{}\allowbreak{}fibre})
\item[Chapter 11] The Lefschetz–Verdier formula, in Ⅱ Duality and the trace formula (\bpslug{chap-{}\allowbreak{}The-{}\allowbreak{}Lefschetz-{}\allowbreak{}Verdier-{}\allowbreak{}formula})
\item[Lemma 11.7] External duality, in Ⅱ Duality and the trace formula (\bpslug{lemma-{}\allowbreak{}external-{}\allowbreak{}duality-{}\allowbreak{}trace-{}\allowbreak{}morphism})
\item[Definition 11.8] Trace morphism, in Ⅱ Duality and the trace formula (\bpslug{definition-{}\allowbreak{}trace-{}\allowbreak{}morphism})
\item[Definition 11.9] Trace class and local terms, in Ⅱ Duality and the trace formula (\bpslug{definition-{}\allowbreak{}trace-{}\allowbreak{}class-{}\allowbreak{}local-{}\allowbreak{}terms})
\item[Theorem 11.13] Proper pushforward and Lefschetz–Verdier, in Ⅱ Duality and the trace formula (\bpslug{theorem-{}\allowbreak{}proper-{}\allowbreak{}pushforward-{}\allowbreak{}lefschetz-{}\allowbreak{}verdier})
\item[Lemma 11.17] Lifting a correspondence to an invariant decomposition, in Ⅱ Duality and the trace formula (\bpslug{lemma-{}\allowbreak{}lift-{}\allowbreak{}correspondence-{}\allowbreak{}invariant-{}\allowbreak{}decomposition})
\item[Lemma 11.18] The closed piece as a pushforward, in Ⅱ Duality and the trace formula (\bpslug{lemma-{}\allowbreak{}closed-{}\allowbreak{}piece-{}\allowbreak{}pushforward})
\item[Chapter 12] The trace formula, in Ⅱ Duality and the trace formula (\bpslug{chap-{}\allowbreak{}The-{}\allowbreak{}trace-{}\allowbreak{}formula})
\item[Remark 12.4] Local terms are not computed stalk by stalk, in Ⅱ Duality and the trace formula (\bpslug{remark-{}\allowbreak{}locality-{}\allowbreak{}warning})
\item[Proposition 12.16] The Frobenius local term on a curve, in Ⅱ Duality and the trace formula (\bpslug{proposition-{}\allowbreak{}frobenius-{}\allowbreak{}local-{}\allowbreak{}term-{}\allowbreak{}curve})
\item[Theorem 12.24] Grothendieck–Lefschetz trace formula, in Ⅱ Duality and the trace formula (\bpslug{theorem-{}\allowbreak{}grothendieck-{}\allowbreak{}lefschetz-{}\allowbreak{}trace-{}\allowbreak{}formula})
\end{description}
\clearpage\phantomsection%
\section*{References}
\addcontentsline{toc}{section}{References}
\bpbibitem{BBD82}{bp102}{A. Beilinson, J. Bernstein, and P. Deligne. \emph{Faisceaux pervers}, volume 100 of \emph{Astérisque}. 1982.}

\bpbibitem{BS15}{bp103}{B. Bhatt and P. Scholze. The pro-étale topology for schemes. \emph{Astérisque}, 369:99–201, 2015. \href{https://people.mpim-bonn.mpg.de/scholze/proetale.pdf}{Author's manuscript}.}

\bpbibitem{Con07}{bp104}{B. Conrad. Deligne's notes on Nagata compactifications. \emph{Journal of the Ramanujan Mathematical Society}, 22:205–257, 2007. \href{https://math.stanford.edu/\~conrad/papers/nagatafinal.pdf}{Author's manuscript}.}

\bpbibitem{Del74}{bp105}{P. Deligne. La conjecture de Weil. I. \emph{Publ. Math. IHÉS}, 43:273–307, 1974.}

\bpbibitem{Del80}{bp106}{P. Deligne. La conjecture de Weil. II. \emph{Publ. Math. IHÉS}, 52:137–252, 1980.}

\bpbibitem{Den22}{bp107}{C. Deninger. There is no “Weil-”cohomology theory with real coefficients for arithmetic curves. \emph{Ann. Sc. Norm. Super. Pisa Cl. Sci.}, 2022. To appear. arXiv:\href{https://arxiv.org/abs/2204.02714}{2204.02714}.}

\bpbibitem{Deu41}{bp108}{M. Deuring. Die Typen der Multiplikatorenringe elliptischer Funktionenkörper. \emph{Abh. Math. Sem. Univ. Hamburg}, 14:197–272, 1941.}

\bpbibitem{Eke90}{bp109}{T. Ekedahl. On the adic formalism. In \emph{The Grothendieck Festschrift, Vol. II}, volume 87 of \emph{Progress in Mathematics}, pages 197–218. Birkhäuser, 1990.}

\bpbibitem{Has36}{bp110}{H. Hasse. Zur Theorie der abstrakten elliptischen Funktionenkörper. I, II, III. \emph{J. Reine Angew. Math.}, 175:55–62, 69–88, 193–208, 1936.}

\bpbibitem{Hum72}{bp111}{J. E. Humphreys. \emph{Introduction to Lie Algebras and Representation Theory}, volume 9 of \emph{Graduate Texts in Mathematics}. Springer, 1972.}

\bpbibitem{Ill94}{bp99}{L. Illusie. Autour du théorème de monodromie locale. \emph{Astérisque}, 223:9–57, 1994. §4. \href{https://www.numdam.org/item/AST_1994__223__9_0/}{Original article}.}

\bpbibitem{Kat88}{bp112}{N. M. Katz. Travaux de Laumon. \emph{Astérisque}, 161–162:105–132, 1988. Séminaire Bourbaki, exp. 691.}

\bpbibitem{KW01}{bp113}{R. Kiehl and R. Weissauer. \emph{Weil conjectures, perverse sheaves and \(\ell\)-adic Fourier transform}, volume 42 of \emph{Ergebnisse der Mathematik}. Springer, 2001.}

\bpbibitem{Lau87}{bp114}{G. Laumon. Transformation de Fourier, constantes d'équations fonctionnelles et conjecture de Weil. \emph{Publ. Math. IHÉS}, 65:131–210, 1987.}

\bpbibitem{Mil86a}{bp115}{J. S. Milne. \href{https://www.jmilne.org/math/xnotes/AVs.pdf}{Abelian Varieties}. In \emph{Arithmetic Geometry}. 1986. Corrected online version.}

\bpbibitem{Mil86b}{bp116}{J. S. Milne. \href{https://www.jmilne.org/math/xnotes/JVs.pdf}{Jacobian Varieties}. In \emph{Arithmetic Geometry}. 1986. Corrected online version.}

\bpbibitem{Ran39}{bp117}{R. A. Rankin. Contributions to the theory of Ramanujan's function \(\tau(n)\) and similar arithmetical functions. II. \emph{Proc. Cambridge Philos. Soc.}, 35:357–372, 1939.}

\bpbibitem{SGA4}{bp118}{M. Artin, A. Grothendieck, and J.-L. Verdier, editors. \emph{Théorie des topos et cohomologie étale des schémas}, volumes 269, 270, 305 of \emph{Lecture Notes in Mathematics}. Springer, 1972–1973. SGA 4.}

\bpbibitem{SGA4½}{bp101}{P. Deligne. \emph{Cohomologie étale}, volume 569 of \emph{Lecture Notes in Mathematics}. Springer, 1977. SGA 4½; \emph{Théorèmes de finitude}, §3, and \emph{Rapport sur la formule des traces}, §3, and \emph{Sommes trigonométriques}, §1.}

\bpbibitem{SGA5}{bp119}{A. Grothendieck, editor. \emph{Cohomologie \(\ell\)-adique et fonctions \(L\)}, volume 589 of \emph{Lecture Notes in Mathematics}. Springer, 1977. SGA 5.}

\bpbibitem{SGA7II}{bp100}{P. Deligne and N. Katz, editors. \emph{Groupes de monodromie en géométrie algébrique}, volume 340 of \emph{Lecture Notes in Mathematics}. Springer, 1973. SGA 7 II; Exposé XIII, (2.1.7).}

\bpbibitem{Stacks}{bp120}{The Stacks Project Authors. \href{https://stacks.math.columbia.edu}{The Stacks Project}. Individual results are cited by their stable tags in the text.}

\bpbibitem{Var05}{bp98}{Y. Varshavsky. Lefschetz–Verdier trace formula and a generalization of a theorem of Fujiwara. 2005. Version 2. arXiv:\href{https://arxiv.org/abs/math/0505564}{math/0505564}. \href{https://arxiv.org/abs/math/0505564v2}{Author's preprint}.}

\bpbibitem{Wei48a}{bp121}{A. Weil. \emph{Sur les courbes algébriques et les variétés qui s'en déduisent}. Hermann, Paris, 1948.}

\bpbibitem{Wei48b}{bp122}{A. Weil. \emph{Variétés abéliennes et courbes algébriques}. Hermann, Paris, 1948.}

\bpbibitem{Wei49}{bp123}{A. Weil. Numbers of solutions of equations in finite fields. \emph{Bull. Amer. Math. Soc.}, 55:497–508, 1949.}

\end{document}
